% BEGIN EMBEDDED BIBLIOGRAPHY
% Bibliography data are part of this source; the compiler writes its temporary input.
\begin{filecontents*}[overwrite]{\jobname.bib}
@unpublished{QED3Master,
 title={The field-strength {Osterwalder--Schrader} axioms for {QED} in three dimensions},
 shorthand={QED}, date={2026}, note = {Comprehensive edition, 263 pp., 20 October 2026, before the modular integration revision},
 url={https://github.com/ethanwsussman/Yang--Mills-LLM-experiment/blob/main/QED3/qed3.pdf}
}
@unpublished{MaxwellBlockManuscript,
 title={Maxwell block minimization and Gaussian fluctuations},
 shorthand={MB}, date={2026}, note={Companion manuscript, version of 20 October 2026},
 url={https://github.com/ethanwsussman/Yang--Mills-LLM-experiment/blob/main/QED3/Submanuscripts/Maxwell%20block%20analysis%20%28MB1%29/MBM1.pdf}
}
@unpublished{FermionBlockManuscript,
 title={Chamber fermion locality and local integration},
 shorthand={FB}, date={2026}, note={Companion manuscript, version of 20 October 2026},
 url={https://github.com/ethanwsussman/Yang--Mills-LLM-experiment/blob/main/QED3/Submanuscripts/Chamber%20fermion%20block%20analysis%20%28FB1%29/FBM1.pdf}
}
@unpublished{NormalizedContraction,
 title={Relevant extraction and contraction of local interactions},
 shorthand={NC}, date={2026}, note={Companion manuscript, version of 20 October 2026},
 url={https://github.com/ethanwsussman/Yang--Mills-LLM-experiment/blob/main/QED3/Submanuscripts/Relevant%20extraction%20and%20normalized%20contraction%20%28RN1%29/RNM1.pdf}
}

@unpublished{NonlinearStep,
 title={The exact nonlinear renormalization step for lattice {QED} in three dimensions},
 shorthand={NL}, date={2026}, note={Companion manuscript, version of 20 October 2026},
 url={https://github.com/ethanwsussman/Yang--Mills-LLM-experiment/blob/main/QED3/Submanuscripts/Nonlinear%20renormalization%20step%20%28NL1%29/NLM1.pdf}
}

@misc{BrightMetric2023,
 author={Bright, Paige}, title={Introduction to Metric Spaces},
 date={2023-01-24}, note={MIT 18.S190, Lecture 5, Theorem 8},
 url={https://ocw.mit.edu/courses/18-s190-introduction-to-metric-spaces-january-iap-2023/mit18_s190iap23_lec5.pdf}}
\end{filecontents*}
% END EMBEDDED BIBLIOGRAPHY
% CVM1 / d042: proof companion to CVO1.
\documentclass[11pt,letterpaper,reqno]{amsart}
% BEGIN INLINE manuscript-style
\makeatletter
\NeedsTeXFormat{LaTeX2e}[2020/10/01]
% Use with \documentclass[11pt,letterpaper,reqno]{amsart}.
% This file controls presentation only. It supplies no mathematical content.

\RequirePackage[T1]{fontenc}
\RequirePackage{lmodern}
\RequirePackage{microtype}
\RequirePackage[letterpaper,margin=1in,headheight=15pt]{geometry}
\RequirePackage{amsmath,amsthm,amssymb,mathtools}
\RequirePackage{enumitem}
\RequirePackage{perpage}
\RequirePackage{xcolor}
\definecolor{manuscriptred}{rgb}{0.6,0.1,0.1}
\definecolor{manuscriptblue}{rgb}{0.1,0.3,0.5}
\definecolor{manuscriptmagenta}{HTML}{60445C}
\RequirePackage[breaklinks=true,colorlinks=true,
  linkcolor=manuscriptred,citecolor=manuscriptblue,urlcolor=manuscriptmagenta]{hyperref}
\RequirePackage{cleveref}
\hypersetup{pdfauthor={}}

% Retain the AMS paragraph, title, heading and running-head conventions.
% Equations retain one global Arabic sequence, as in the reference.
\setcounter{tocdepth}{1}
\renewcommand{\thefigure}{\Roman{figure}}
\renewcommand{\thetable}{\Roman{table}}
\MakePerPage{footnote}
\renewcommand{\thefootnote}{%
  \ifnum\value{footnote}>9\arabic{footnote}\else\fnsymbol{footnote}\fi}

% Upright statements. Main results are A, B, C; supporting statements
% share a section-local sequence. These end marks indicate typography,
% never proof completeness, validation, or a deferred-proof status.
\theoremstyle{definition}
\newtheorem{ms@theorem}{Theorem}
\renewcommand{\thems@theorem}{\Alph{ms@theorem}}
\newtheorem*{ms@theoremstar}{Theorem}
\newtheorem{ms@proposition}{Proposition}[section]
\newtheorem{ms@lemma}[ms@proposition]{Lemma}
\newtheorem{ms@definition}[ms@proposition]{Definition}
\newtheorem{ms@corollary}[ms@proposition]{Corollary}
\newtheorem*{ms@corollarystar}{Corollary}

\newcommand{\ms@majorstart}{%
  \pushQED{\qed}\renewcommand{\qedsymbol}{$\blacksquare\square$}}
\newcommand{\ms@minorstart}{%
  \pushQED{\qed}\renewcommand{\qedsymbol}{$\blacksquare$}}
\newenvironment{theorem}[1][]{\ms@majorstart\ms@theorem[#1]}
  {\popQED\endms@theorem}
\newenvironment{theorem*}[1][]{\ms@majorstart\ms@theoremstar[#1]}
  {\popQED\endms@theoremstar}
\newenvironment{proposition}[1][]{\ms@minorstart\ms@proposition[#1]}
  {\popQED\endms@proposition}
\newenvironment{lemma}[1][]{\ms@minorstart\ms@lemma[#1]}
  {\popQED\endms@lemma}
\newenvironment{definition}[1][]{\ms@minorstart\ms@definition[#1]}
  {\popQED\endms@definition}
\newenvironment{corollary}[1][]{\ms@majorstart\ms@corollary[#1]}
  {\popQED\endms@corollary}
\newenvironment{corollary*}[1][]{\ms@majorstart\ms@corollarystar[#1]}
  {\popQED\endms@corollarystar}
% The source's lemma/proposition variants are deliberately not carried over:
% choose one marker for each statement type rather than encode proof status.

\theoremstyle{remark}
\newtheorem{remark}[ms@proposition]{Remark}
\newtheorem*{remark*}{Remark}
% The ordinary amsthm proof environment retains its open square.

\crefname{ms@theorem}{theorem}{theorems}
\Crefname{ms@theorem}{Theorem}{Theorems}
\crefname{ms@proposition}{proposition}{propositions}
\Crefname{ms@proposition}{Proposition}{Propositions}
\crefname{ms@lemma}{lemma}{lemmas}
\Crefname{ms@lemma}{Lemma}{Lemmas}
\crefname{ms@definition}{definition}{definitions}
\Crefname{ms@definition}{Definition}{Definitions}
\crefname{ms@corollary}{corollary}{corollaries}
\Crefname{ms@corollary}{Corollary}{Corollaries}
\crefname{remark}{remark}{remarks}
\Crefname{remark}{Remark}{Remarks}

% Bibliography configuration for this manuscript.
  \RequirePackage[backend=biber,style=alphabetic,backref=true,
    giveninits=true,isbn=false,date=year]{biblatex}
  \addbibresource{\jobname.bib}
  \setcounter{biburlnumpenalty}{100}
  \setcounter{biburllcpenalty}{100}
  \setcounter{biburlucpenalty}{100}
  \renewcommand*{\bibfont}{\footnotesize}
  \renewbibmacro{in:}{}
  \DeclareFieldFormat{pages}{#1}
  % Use the public alphabetic style. No private author-hash overrides;
  % repeated author names are printed in full.
\makeatother
% END INLINE manuscript-style
\usepackage{mathrsfs}
\newcommand{\R}{\mathbb R}
\newcommand{\C}{\mathbb C}
\newcommand{\Z}{\mathbb Z}
\newcommand{\N}{\mathbb N}
\newcommand{\Sspace}{\mathcal S}
\newcommand{\dd}{\mathrm d}
\DeclareMathOperator{\supp}{supp}
\DeclareMathOperator{\tr}{tr}
\DeclareMathOperator{\diag}{diag}
\newcommand{\extref}[2]{#1 of~\cite{#2}}
\newcommand{\exteq}[2]{(#1) of~\cite{#2}}
\allowdisplaybreaks
\title[Agreement across covers and common tuning]{Agreement across covers and common tuning\\for lattice QED in three dimensions}
\date{}
\hypersetup{pdftitle={Agreement across covers and common tuning for lattice QED in three dimensions},pdfauthor={}}
\begin{document}
\begin{abstract}
We construct one mass and vacuum pair for the prescribed family of odd
reference tori in three-dimensional chamber lattice QED, under explicit
primitive and nonlinear-step hypotheses. The pair is uniquely selected
among controlled trajectories by zero terminal couplings on the top
torus. Signed projection of the primitive kernels and comparison on
nonwrapping local lifts identify the actual renormalization steps on
smaller odd quotients. A backward mass equation and a forward activity
equation give a finite-sequence contraction; continuation and local
comparison then transfer the selected pair with uniform bounds.
\end{abstract}
\maketitle
\tableofcontents

\section{Introduction}\label{sec:introduction}
Fix \(0<\rho<1\), the chamber weights \((6,6\sqrt2,6\sqrt2,6)\),
and an odd integer \(L\ge4^8\), an integral power \(M\) of \(L\) with
\(M\ge L^2\), integers \(r\ge2\) and \(p_0,p>0\), \(0<\varepsilon\le1/200\),
\(1-\varepsilon<\alpha<1\), \(c_{\rm fl}\ge1\), and \(0<\bar m\le1\).
The photon is periodic and the four-component chamber spinors are
antiperiodic in all three directions. Each photon site has 48 chamber
labels, with measure \(\eta_k^3/48\) at mesh \(\eta_k=L^{-k}\).
For an integer \(N\ge0\), set
\[
 h=L^{\lfloor\log_L(1/\bar m)\rfloor},\qquad
 \eta=L^{-N},\quad K=N+\log_Lh,\quad
 a_j=L^{j-N},\quad e_j=e\sqrt{a_j},\quad \mu_j=\bar m a_j.
\]
Here \(0\le j\le K\), and \(q=e\sqrt\eta\) denotes the charge at
level zero. We write \(\bar e_0,\bar e_1\) for the fixed charge
ceilings, to distinguish them from the running charges \(e_j\).
Put \(v_j=-\log e_j\), and let \(r_j\) be the least power \(L^i\),
\(i\in\mathbb Z_{\ge0}\), with \(r_j\ge v_j^r\). The physical sides are
\[
 R_{\rm top}=hMr_K,\quad R^*=R_{\rm top}/L,\qquad
 \mathcal T(e)=\{hL^T:T\in\mathbb N\cup\{0\},\ hL^T\le R_{\rm top}\}.
\]
They are independent of \(N\). For \(R\in\mathcal T(e)\), define
\(\ell_j=R/a_j\) and
\(\mathfrak n_j=\log_L(\ell_j/(Mr_j))\). The stopping level \(k_*(R)\)
is the first \(j\) with \(\mathfrak n_j=0\), when it exists.
Write \(k_{\rm top}=k_*(R_{\rm top})\). A star denotes evaluation
at \(k_*(R)\): for example \(a_*=a_{k_*(R)}\) and
\(\mu_*=\mu_{k_*(R)}\).

\paragraph{Numerical hypotheses.}
Assume
\[
 \begin{gathered}
 0<e\sqrt h\le\bar e_1\le\min(\bar e_0,1/3),\qquad
 \bar e_0^\varepsilon\le1/4,\\
 c_{\rm fl}L^{-3/8+6/200}\le1/16,\qquad L^{-5/8}\le1/4,
 \end{gathered}
\]
and, for every \(v\ge-\log\bar e_1\),
\begin{equation}\label{eq:cv-scale-conditions}
 \begin{gathered}
 v\ge8,\qquad
 \left(1+\frac{3\log L}{2v}\right)^r<L,\qquad
 v^r>31L+16,\\
 (c_{\rm fl}+2)e^{-\varepsilon v}
       \le\frac{1-L^{-\varepsilon/2}}{16}.
 \end{gathered}
\end{equation}

\paragraph{The activity maps.}
Use the real chamber activity spaces, with independent connection
arguments, defined in \cref{sec:activities}. Write \(\mathscr E_j\) for an activity
with external connection \(\mathscr C\), and \(E_j=\mathscr E_j|_{\mathscr C=0}\).
At every used scale impose all the simultaneous hypotheses
(i)--(iv) of \cite[Theorem~A]{NonlinearStep}, including its
primitive, inverse-domain, determinant, extraction, reconstruction,
contour and connected-logarithm conditions. The constants, norms and
operators are those of that theorem and of
\cite[Theorems~A--B]{NormalizedContraction}; their definitions and the
parts used below are recalled in \cref{sec:activities,sec:inputs}.
These are hypotheses on the actual chamber maps. In particular the
following tuning inequalities do not replace their domain or numerical
conditions. The decay weight satisfies \(\kappa\ge7\kappa_0+7\), where
\(\kappa_0=54\log72\); the constants are calibrated together as in the
cited step theorem.

On the top torus choose the photon and fermion families of
\cref{subsec:rgcov-top-system}, with polymer decay
\(\kappa_w+\kappa_0\), where \(\kappa_w=\kappa+\kappa_0+3\).
On each smaller odd quotient use their projections, with common
constant \(C_w=64K_{\mathrm{pol}}C_w^{(1)}\), where
\(K_{\mathrm{pol}}=\exp(\kappa_0)/36\) is the scalar polymer-sum
constant of \cite[Appendix~B, (81)]{NormalizedContraction}.
Both incoming and outgoing sides are required to have
\(\ell/(LM)\) odd and at least 13. The step comparison requires
\(\ell_{\rm small}/L>(31L+16)M\); the stopping calculation verifies
these side conditions at all steps used in the theorem.

For a comparison with smaller unit side $\ell_1$, a connected
covering polymer $\widetilde X$ is admissible when its smallest
rectangular box has largest side at most $\ell_1-14(2L+1)M$.
Its two torus images are called flat. Local agreement means equality
under the common covering generators for connection pairs agreeing
on $\operatorname{box}(\widetilde X)^{\sim6(2L+1)}$.
Here the superscript denotes enlargement by that number of $M$-cube
layers. \Cref{prop:rgcov-common-lifts} justifies this identification.

The estimates used for tuning take the following form.
Set \(w_j^{\mathrm m}=e_j^{3/4-8\varepsilon}\) and
\(w_j^{\mathrm E}=e_j^{1/4-7\varepsilon}\).
The no-slot data norm is
\(\max\{|m|/w_j^{\mathrm m},\|E\|_j/w_j^{\mathrm E}\}\); with an external slot,
replace \(\|E\|_j\) by \(\|\mathscr E\|_j^x\).
Radius-four data balls mean the corresponding closed balls in the
complexified spaces. Unit balls below are closed balls in their real
subspaces. At zero the activity coordinate is fixed to zero.
The maps preserve reality and obey
\begin{align}
 m_{j+1}&=L(m_j+m(E_j)),&
 \varepsilon_{j+1}^v&=L^3(\varepsilon_j^v+\varepsilon_v(E_j)+\delta\varepsilon_j^0),
 \label{eq:cv-scalar-recursion}\\
 E_{j+1}&=\Phi_j(m_j,E_j),&
 \Phi_j(m,E)&=\mathcal L(RE)^{\rm nat}
                  +\mathcal LE_j^{\det}+E_j^*(m,E).
 \label{eq:cv-activity-recursion}
\end{align}
Restriction to \(\mathscr C=0\) commutes with the full external-slot recursion
and does not increase its norm. Scalar extraction is linear and real,
\[
 |m(E)|\le c_{\rm fl}e_j^{1/2}\|E\|_j,\qquad
 |\varepsilon_v(E)|\le c_{\rm fl}\|E\|_j,\qquad
 0\le\delta\varepsilon_j^0\le e_j^7.
\]
On the unit ball the normalized size and difference estimates are
\begin{equation}\label{eq:cv-normalized-inputs}
 \begin{aligned}
 L^{-1}w_{j+1}^m/w_j^{\mathrm m}&\le1/4,\\
 \|\Phi_j(m,E)\|_{j+1}/w_{j+1}^E
   &\le\tfrac1{16}\|E\|_j/w_j^{\mathrm E}+2e_{j+1}^{\varepsilon},\\
 |m(E)-m(\widetilde E)|/w_j^{\mathrm m}
   &\le c_{\rm fl}e_j^\varepsilon d_E,\\
 \|\Phi_j(m,E)-\Phi_j(\widetilde m,\widetilde E)\|_{j+1}/w_{j+1}^E
   &\le\tfrac1{16}d_E+e_{j+1}^\varepsilon(d_m+d_E),
 \end{aligned}
\end{equation}
where \(d_m=|m-\widetilde m|/w_j^{\mathrm m}\) and
\(d_E=\|E-\widetilde E\|_j/w_j^{\mathrm E}\).
The same size estimate holds for the external-slot norm. The
fixed-point argument uses the no-slot difference estimate.
Initially \(E_0=\mathscr E_0=\mathfrak P_0=0\), where
\(\mathfrak P_j\) is the scalar external family in \cref{sec:inputs}.
The first step uses the chamber definitions in \cref{sec:inputs}.

\begin{theorem}[A common mass and vacuum pair]\label{thm:main}
Under the hypotheses above, the following assertions hold.
For each \(R\in\mathcal T(e)\) there is a finite threshold
\(N_{\min}(R,e)\) such that the first stop exists for
\(N\ge N_{\min}(R,e)\). Its physical quantities are then independent
of \(N\). The thresholds are nonincreasing and the stopping levels
nondecreasing
with \(R\); their precise formulas are in \eqref{eq:rgfl-mesh-threshold}.

For \(N\ge N_{\min}(R_{\rm top},e)\), there is exactly one real
pair \(P_N(e)=(m_{0,N}(e),\varepsilon_{0,N}^v(e))\) whose top-torus
trajectory is controlled and has
\(m_{k_{\rm top}}=\varepsilon_{k_{\rm top}}^v=0\).
Controlled means that each preceding step is defined in its
radius-four data ball and \(|m_j|\le w_j^{\mathrm m}\) through the stop.
For every \(R\in\mathcal T(e)\) with \(N\ge N_{\min}(R,e)\),
the same pair is controlled on \(R\). Up to its stop, the scalar
mass and vacuum sequences equal the top sequences, while the
activities and external scalar families agree on admissible flat
lifts under the algebra identification of \cref{prop:rgcov-common-lifts}. On each such torus,
\begin{equation}\label{eq:cv-common-bounds}
 |m_j|\le\tfrac18e_j^{3/4-8\varepsilon},\qquad
 \|E_j\|_j\le\|\mathscr E_j\|_j^x\le4e_j^{1/4-6\varepsilon},\qquad
 |\varepsilon_j^v|\le5c_{\rm fl}e_j^{1/4-6\varepsilon}.
\end{equation}
The smaller terminal couplings need not vanish.

Within the fixed construction, the pair is determined by
\((N,e,\bar m)\) and is independent of the smaller side, twist,
external connection and probe. Define only \(m_{0,N}(e)=0\) below
the top definition threshold. For every \(N\ge0\),
\[
 m_{0,N}(e)\in\mathbb R,\qquad
 |m_{0,N}(e)|\le\tfrac18q^{3/4-8\varepsilon}<\tfrac18 .
\]
The total bare mass \(\boldsymbol\mu_N=\bar mL^{-N}+m_{0,N}(e)\)
has no asserted sign. On \(R^*\), for
\(N\ge N_{\rm ref}(e):=N_{\min}(R^*,e)\),
\[
 a_*\in\{h/L,h/L^2\},\qquad
 \ell_*=Mr_*,\qquad r_*>31L+16,\qquad
 L^{-3}<\mu_*\le L^{-1}.
\]


\end{theorem}

The theorem concerns the prescribed family of odd reference tori.
The scalar pair is selected before inserting a smaller side, a harmonic
twist or an external probe. The proof has two parts. On the top torus,
the mass equation is solved backwards from zero terminal mass, while
the activity equation runs forwards from zero initial activity. In the
weighted sequence norm these two equations define a contraction with
factor \(5/16\). The vacuum is recovered by an explicit backward sum.

To use that pair on a smaller torus, we compare the actual local
operations. Odd covers preserve the antiperiodic character. Projection
gives a signed finite image sum; when weakening parameters confine a
row to a common nonwrapping lift, only one image survives. The local
Grassmann algebra then identifies extraction, determinant traces, both
finite localizations and the fluctuation integrals. Absolute convergence
passes this equality to the connected logarithm. At each smaller-volume
step, its own continuation estimate controls the full activity norm,
and local comparison identifies the scalar couplings with those of the
top trajectory.

The cover and tuning arguments are developed in
\cite[\S\S15--16]{QED3Master}.
The photon, fermion, extraction and nonlinear constructions are supplied
by \cite{MaxwellBlockManuscript,FermionBlockManuscript,NormalizedContraction,NonlinearStep}.
We give the cover comparison and finite-sequence tuning proofs here,
with those constructions as explicit inputs.

The definitions in \cref{sec:activities} and the input identities in
\cref{sec:inputs} fix the maps compared in
\cref{sec:rg-cover-comparison,sec:local-comparison}.
The stopping calculation is in \cref{sec:rg-flow};
\cref{sec:common-tuning} proves the common-pair theorem.
\section{Activities, connections and coefficient norms}
\label{sec:activities}

We specify the local activity space used in \cref{thm:main}.
These are the conventions of \cite[\S2]{NormalizedContraction}.

\subsection{Chamber sites, pairs and connection domains}

Write \(O_{\mathrm{cub}}\) for the signed permutation matrices on
\(\R^3\). A chamber site is \(x=(c,g)\), with photon projection
\(\pi x=c\) and \(g\in O_{\mathrm{cub}}\). Integration over chamber
sites uses \(\dd x=(\eta_k^3/48)\sum_{c,g}\). We take
\(M\)-cubes centered at \(M\Z^3/\ell\Z^3\); their lattice sites
form a partition. Blocks of every odd side used below are centered
in the same way.

For the three internal neighbors let \(s_0\) change the sign of
the third coordinate, let \(s_1\) exchange the second and third
coordinates, and let \(s_2\) exchange the first and second.
Let \(r_{e_1}\) change the sign of the first coordinate. Set
\[
 x^{(i)}=(c,gs_i)\quad(0\le i\le2),\qquad
 x^{(3)}=(c+\eta_kge_1,gr_{e_1}).
\]
Each neighbor map is an involution. An oriented bond field
\(B\) is antisymmetric under reversal and periodic on the photon
lattice. Write \(B_\mu(c)=B(c,c+\eta_ke_\mu)\), and use the
forward differences
\[
 \partial_\nu B_\mu(c)=
       \eta_k^{-1}\bigl(B_\mu(c+\eta_ke_\nu)-B_\mu(c)\bigr),
 \qquad
 (d\lambda)_\mu(c)=
       \eta_k^{-1}\bigl(\lambda(c+\eta_ke_\mu)-\lambda(c)\bigr).
\]
Thus \((dB)_{\nu\mu}=\partial_\nu B_\mu-\partial_\mu B_\nu\).
Distances between photon sites are sup distances, using the
shortest lift when a distance bound is imposed. For \(0<\alpha<1\)
the three norms are
\[
 \|B\|_1=\|B\|_\infty,\qquad
 \|B\|_2=\|\partial B\|_\infty,\qquad
 \|B\|_3=\|\delta_\alpha\partial B\|_\infty
 =\sup_{\substack{\mu,\nu,c,c'\\0<d(c,c')\le1}}
 \frac{|\partial_\nu B_\mu(c')-\partial_\nu B_\mu(c)|}
      {d(c,c')^\alpha}.
\]
Local versions use only the indicated cells in the given region.
Let \(\square^\dagger\) be the concentric enlargement of an
\(M\)-cube to side \((2L+1)M\). Put
\(\theta_i=e_k^{-3/4+i\varepsilon}\).
The connection domains are
\begin{equation}\label{eq:connection-domains}
\begin{split}
 R_k&=\{B:\|B\|_i<\theta_i,\ 1\le i\le3\},\\
 \widetilde R_k&=\{B_0+B_1:B_1\in R_k,\ B_0\text{ real, and on each }
          \square^\dagger,\\
 &\hspace{25mm} B_0=B_\square^{\mathrm b}+d\lambda_\square
   \text{ with }B_\square^{\mathrm b},\lambda_\square\text{ real and }
          \|B_\square^{\mathrm b}\|_i<\theta_i\},\\
 \mathfrak X_k&=\{\mathscr C:\|\mathscr C\|_\infty<e_k^{-1+2\varepsilon},\
                  \|\partial\mathscr C\|_\infty<e_k^{-1+3\varepsilon}\}.
\end{split}
\end{equation}
Only existence of these representatives is required.
For \(0<c\le1\), the notation \(c\widetilde R_k\) multiplies
all three bounds for both summands by \(c\).
For \(\widetilde R_k^-\) the local real bounds remain
\(\theta_i\), while the global summand has bounds
\((1-e_k^\varepsilon)\theta_i\).

For a short ordered pair \(x=(c,g),x'=(c',g')\), let
\(\tau B(c,c')\) be the arithmetic mean of the sums of \(B\)
along the six coordinate-order paths from \(c\) to \(c'\).
The bond values in this sum are unscaled; define
\[
 U_B(x,x')=\exp\{ie_k\eta_k\tau B(c,c')\},\qquad
 (\delta_{\alpha,B}\psi)(x,x')
       =\frac{U_B(x,x')\psi(x')-\psi(x)}
                    {\max(|c-c'|_\infty,\eta_k)^\alpha}.
\]
For pairs within one cluster the transport is one. Barred fields
use \(U_B^{-1}\). These formulas specify evaluation; the pair
generators in the algebra remain independent.

We use Pauli matrices and
\[
 \gamma_\mu=\sigma_3\otimes\sigma_\mu,\quad
 \gamma_5=\sigma_1\otimes1_2,\quad \Sigma_3=\sigma_3\otimes1_2,
 \qquad
 \mathsf C=i\gamma_2,\quad
 \varkappa=\gamma_5\mathsf C=\sigma_2\otimes\sigma_2 .
\]
Here \(\mathsf C\) is real with
\(\mathsf C^T=\mathsf C^{-1}=-\mathsf C\), while
\(\varkappa\) is a real symmetric involution. The spin group below
is the group generated by \(\gamma_1,\gamma_2,\gamma_3,\gamma_5\).
Its diagonal coordinate action \(r(S)\) is determined by
\(S\gamma_\mu S^{-1}=r(S)_{\mu\mu}\gamma_\mu\).
Connections transform by pullback on oriented bonds.

\subsection{The local Grassmann algebra}
\label{subsec:rg-activity-spaces}

Fix a level $k$, mesh $\eta_k=L^{-k}$ and an admissible side
$\ell$, with $\ell/(LM)$ odd and at least $13$.
Let $\mathcal P_k$ be the collection of face-connected polymers
of $M$-cubes on this torus. Write $|X|_M$ for the number of its $M$-cubes. The number $Md_M(X)$
is the infimum of the total Euclidean length of a connected finite
net of straight segments meeting every closed cube of $X$; the
infimum includes the segment representations on the torus.
The notation $X^{\sim n}$ means enlargement by $n$ layers of
$M$-cubes in the sup metric.
A polymer is small when
$d_M(X)\le L$. A small polymer has an injective lift, unique up
to a period translation, by the diameter and lifting estimates
of \cite[Appendix~B]{NormalizedContraction}.
Write $\operatorname{Vol}(X)=M^3|X|_M$.

The independent Grassmann generators are $\psi$, $\bar\psi$ on
the chamber sites $x=(c,g)$ and $\chi$, $\bar\chi$ on ordered
pairs. For distinct chamber sites on the covering lattice, put
$d_X(x,x')=\max(|c-c'|_\infty,\eta_k)$ and restrict pairs to
$d_X\le M$. Pairs on the torus are simultaneous period classes;
their shortest lifts are unique since $2M<\ell$. A pair is
admissible for $X$ when its photon sites and all six rectilinear
paths between them lie in $X$, or when it is $(x,x^{(3)})$ and
one of its photon sites belongs to $X$.
Thus a pair may leave $X$ only along a bond of type $3$.
Pairs inside a cluster have $d_X=\eta_k$ and are admissible when
their common photon site belongs to $X$.
Ordinary generators acquire the sign $(-1)^{n_1+n_2+n_3}$
under translation by $\ell n$. Pair generators acquire the same
sign when both endpoints are translated by $\ell n$.
The kernels are antiperiodic in each ordinary argument and in
each simultaneous pair argument with these conventions.

Let $\xi=(x,\beta,\omega)$ index an ordinary leg and
$\zeta=(x,x',\beta,\omega)$ a H\"older leg, where
$\omega=0$ is unbarred and $\omega=1$ is barred.
We use the index measures
\[
 \nu(\xi)=\eta_k^3/48,\qquad
 \nu(\zeta)=(\eta_k^3/48)^2.
\]
Integrals sum the spin and bar indices as well. An activity
has a unique expansion with separately antisymmetric kernels,
\begin{equation}\label{eq:rg-activity-kernels}
 \mathscr E(X;\mathscr A,\mathscr C)
 =\sum_{n,m}\frac1{n!m!}\int\dd\xi\,\dd\zeta\,
      \mathscr E_{nm}(X;\mathscr A,\mathscr C;\xi;\zeta)
                       \psi(\xi_1)\cdots\psi(\xi_n)
                       \chi(\zeta_1)\cdots\chi(\zeta_m).
\end{equation}
Barred generators are included in this notation.
The generator weights, as distinct from the index measures, are
\[
 h_{k,1}=e_k^{-1/4},\qquad h_{k,2}=e_k^{-1/4+\varepsilon}.
\]
For an element $F$ the weighted coefficient norm is
\begin{equation}\label{eq:rg-activity-coefficient-norm}
 \|F\|_{h_k}
       =\sum_{n,m}\frac{h_{k,1}^nh_{k,2}^m}{n!m!}
                                   \int\dd\xi\,\dd\zeta\,|F_{nm}|.
\end{equation}
In particular a generator has norm $h_{k,i}$;
the weight $\nu$ enters only when coefficients are written as
integral kernels.

An activity family in $K_k^{x,\C}$ is balanced: every monomial
has as many barred as unbarred legs. Its kernels are holomorphic
on $\widetilde R_k\times\mathfrak X_k$, vanish unless all ordinary
legs lie over $X$ and all pairs are admissible for $X$, and read
the two connections only on bonds meeting $X$. The norm is
\begin{equation}\label{eq:rg-activity-norm}
 \|\mathscr E\|_k^x
  =\sup_{X\in\mathcal P_k}\exp(\kappa d_M(X))
    \sup_{(\mathscr A,\mathscr C)\in\widetilde R_k\times\mathfrak X_k}
                          \|\mathscr E(X;\mathscr A,\mathscr C)\|_{h_k}.
\end{equation}
We require $\kappa\ge\kappa_0$ and finite norm.

The required symmetries act on both the connections and the
generators. For $\tilde g=(S,a)$ with the coordinate sign action
$r(S)$ and $a\in M\Z^3$, put
\[
 \tilde g(c,g)=(r(S)c+a,r(S)g),\qquad
 a_{\tilde g}\psi(x)=S\psi(\tilde g^{-1}x),\qquad
 a_{\tilde g}\bar\psi(x)=(S^{-1})^T\bar\psi(\tilde g^{-1}x).
\]
The same spin matrices act on pair generators after transforming
both endpoints. Local gauge transformations multiply an unbarred
ordinary or pair leg by $\exp(ie_k\Lambda(c))$, using the photon
site of its first point, and a barred leg by the inverse phase.
Write this automorphism as $G_\Lambda$.
Charge conjugation and the antilinear reality map act by
\[
 \mathrm{Cg}\,\psi=\mathsf C\bar\psi,\qquad
 \mathrm{Cg}\,\bar\psi=-\mathsf C\psi,\qquad
 \mathbf K\psi=\varkappa\psi,\qquad
 \mathbf K\bar\psi=\varkappa\bar\psi,
\]
and by the same rules on the H\"older generators.

The covariance requirements on a family are
\begin{equation}\label{eq:rg-activity-symmetries}
 \begin{split}
 a_{\tilde g}\mathscr E(\tilde gX;\tilde g\mathscr A,
                                      \tilde g\mathscr C)
                                  &=\mathscr E(X;\mathscr A,\mathscr C),\\
 \mathrm{Cg}\,\mathscr E(X;-\mathscr A,-\mathscr C)
                                  &=\mathscr E(X;\mathscr A,\mathscr C).
 \end{split}
\end{equation}
For real scalars $\lambda,\lambda'$ and two pairs in the domain
which satisfy $\mathscr A^1=\mathscr A-d\lambda$ and
$\mathscr C^1=\mathscr C-d\lambda'$ on bonds at distance at
most one from $X$, the local gauge requirement is
\[
 G_{\lambda+\lambda'}\mathscr E(X;\mathscr A^1,\mathscr C^1)
                               =\mathscr E(X;\mathscr A,\mathscr C).
\]
The real subspace $K_k^x$ has in addition
\[
 \mathbf K\mathscr E(X;-\overline{\mathscr A},-\overline{\mathscr C})
                               =\mathscr E(X;\mathscr A,\mathscr C).
\]
All requirements refer to the generators before evaluating fields.
They are essential to the form of the counterterms.

The spaces $K_k^\C,K_k$ without the second connection argument
have the same definitions on $\widetilde R_k$. Restriction
$\mathscr E|_0(X;\mathscr A)=\mathscr E(X;\mathscr A,0)$ maps
$K_k^{x,\C}$ into $K_k^\C$ and has norm at most one. The real
subspaces also restrict. Analyticity, locality and the symmetries
follow by setting the second argument and its gauge function to
zero; the norm bound follows from the smaller supremum.

\subsection{Total connections}
\label{subsec:rg-total-connection-domain}

The comparison domain $\widetilde G_k=\widetilde{\mathcal G}_k$ consists of fields $B_0+B_1$
with $B_0$ real and $\|B_1\|_\infty<e_k^{-1+\varepsilon}$,
where on each $\square^{\sim5}$ the field $B_0$ has a real-gauge
representative of supremum norm below $e_k^{-1+\varepsilon}$.
Here $\square^{\sim5}$ has the same center and side $11M$.
The real-gauge representatives in the domains defined above may
be chosen independently on each enlarged cube.

Assume $L\ge5$, $0<\varepsilon\le1/200$ and
\begin{equation}\label{eq:rg-domain-charge-conditions}
 e_k^\varepsilon\le\tfrac14,\qquad
 e_k^{1/4}+e_k^\varepsilon\le1,\qquad
 4e_k^{1/4+2\varepsilon}+2e_k^{3\varepsilon}\le1.
\end{equation}
Then
\begin{equation}\label{eq:rg-total-domain-inclusion}
 \widetilde R_k+\mathfrak X_k\subset\widetilde G_k.
\end{equation}
Indeed, write $\mathscr A=A_0+A_1$ with $A_1\in R_k$ and
$A_0$ locally real-gauge equivalent on $\square^\dagger$ to
$A_\square^b$ with the three threshold bounds. Restrict the
representative to $\square^{\sim5}\subset\square^\dagger$.
Its supremum norm is below
$e_k^{-3/4+\varepsilon}\le e_k^{-1+\varepsilon}$, and
\[
 \|A_1+\mathscr C\|_\infty
 <e_k^{-3/4+\varepsilon}+e_k^{-1+2\varepsilon}
 =(e_k^{1/4}+e_k^\varepsilon)e_k^{-1+\varepsilon}
 \le e_k^{-1+\varepsilon}.
\]
This is the required decomposition of the total connection.
Real gauge transformations can absorb any real constant part,
locally on the nonwrapping cubes.

Put $\mathfrak i(B)=e_k\|\operatorname{Im}B\|_\infty$ and
$q_d(B)=e_k\|dB\|_\infty$. The same decomposition gives
\begin{equation}\label{eq:rg-total-imaginary-curvature-bounds}
 \mathfrak i(\mathscr A+\mathscr C)
       \le e_k^{1/4+\varepsilon}+e_k^{2\varepsilon}\le1,
 \qquad
 q_d(\mathscr A+\mathscr C)
       \le4e_k^{1/4+2\varepsilon}+2e_k^{3\varepsilon}\le1.
\end{equation}
For the curvature, every plaquette lies in an enlarged cube,
and $dA_0=dA_\square^b$ there. This costs at most
$2\|\partial A_\square^b\|_\infty$; the global small field
costs another $2\|\partial A_1\|_\infty$, and the external
field costs $2\|\partial\mathscr C\|_\infty$.
The imaginary bound uses only that $A_0$ is real.

The domains are open: the strict inequalities define open
sets for the global summand, which can absorb a sufficiently
small complex perturbation while the local real representative
is kept fixed. The norm-defined sets are invariant under sign
and conjugation; the local real-gauge definitions are also
invariant under real gauges and real constant additions.
Coordinate sign changes and $M$ translations preserve them.
Their derivative and H\"older norms are preserved, with a
possible reversal of a bond and translation of its base point.
Thus \eqref{eq:rg-total-domain-inclusion} applies also at the
signs, conjugates and real gauges used below.



\subsection{The operator and photon data}

The injection and cluster mean are
$(Jv)(c,g)=v(c)$ and $(J^*f)(c)=48^{-1}\sum_gf(c,g)$.
The normalized average $\mathcal Q_k=Q_kJ^*$, its reverse
$\mathcal Q_k^\sharp=\mathcal Q_k(-B)^T$, the one-step average
$Q$, and the chamber operator $D_B$ are given in
\cite[\S1]{FermionBlockManuscript}. Put $b_k=(1-L^{-1})/(1-L^{-k})$ for $k\ge1$. At these levels set
\[
 S_k(B)=(D_B+\mu+b_k\mathcal Q_k^\sharp\mathcal Q_k)^{-1},\quad
 H_k(B)=b_kS_k(B)\mathcal Q_k^\sharp(B),\quad
 D_k(B)=b_k-b_k^2\mathcal Q_k(B)S_k(B)\mathcal Q_k^\sharp(B).
\]
The one-step projection is $P_L(B)=Q^T(-B)Q(B)$. The barred map is $\bar H_k(B)=\mathsf C^{-1}H_k(-B)\mathsf C$.
At $k=0$, $H_0=\bar H_0=1$ and $D_0(B)=D_B+\mu$.
The symbol $\Sspace(a,\ell)$ denotes the four-component spinor
space on the $a$-spaced torus, with measure $a^3$; the superscript
$\mathrm{ch}$ adds chamber labels and divides the measure by $48$.
All spaces have the antiperiodic convention stated above.
For later use, the hopping notation of \cite[\S2.1]{FermionBlockManuscript} is
\[
 R_i(B)f(x)=U_B(x,x^{(i)})f(x^{(i)}),\qquad
 \mathcal L_B=\sum_{i=0}^3w_i(1-R_i(B)),\qquad
 M_i(c,g)=\tfrac{w_i}{2}(\rho+\gamma\cdot g\mathbf n_i),
\]
where $(w_0,w_1,w_2,w_3)=(6,6\sqrt2,6\sqrt2,6)$ and
$(\mathbf n_0,\mathbf n_1,\mathbf n_2,\mathbf n_3)=(e_3,(e_2-e_3)/\sqrt2,(e_1-e_2)/\sqrt2,-e_1)$.
Thus $\widehat d_i(c,g)=-g\mathbf n_i$. On a nearest-neighbor photon bond
write $u_B(c,c')=\exp(ie_k\eta_kB(c,c'))$; this is the transport
$U_B(x,x^{(3)})$ when $c,c'$ are its endpoints.

For photon fields use the Landau reconstruction $H_k^{\mathrm{ph}}$,
the constrained minimum $A_k^{\min}$ and the real coordinate map
$C:V_k\to\mathcal Z$ of
\cite[Theorem~A]{MaxwellBlockManuscript}. Write
$C_k^{1/2}=\widetilde\Delta_k^{-1/2}$ and
$\mathrm{Sq}_k=CC_k^{1/2}$. The coordinate basis $\{\Upsilon\}$ is
orthonormal and has the specified block or face-bond supports.
The law $\mu_I$ is the product standard Gaussian in these coordinates.
For the fixed exponents $p_0,p$ write
$p_{0,k}=(-\log e_k)^{p_0}$ and $p_k=(-\log e_k)^p$.
Here $\mathrm{Box}_k=\{\sup_\Upsilon|\zeta(\Upsilon)|\le p_{0,k}\}$;
its restriction to $\mathrm{Box}_k$ has probability
$\exp(-\ell^3\delta\varepsilon_k^0)$, with
$0\le\delta\varepsilon_k^0\le12\exp(-p_{0,k}^2/2)$.
The symbols $C$ and $\mathscr C$ always denote, respectively,
this real coordinate map and the external connection.

The common constant $c_\gamma$ bounds both
\[
 \sup_\Upsilon\sum_{\Upsilon'}\exp[-\gamma_wd(\Upsilon,\Upsilon')],
 \qquad \sup_b\sum_{b'}\exp[-\gamma_wd(b,b')],
\]
where the second sum is over unit bonds. The constants $C_w$ and
$\gamma_w$ dominate the photon bounds and the corresponding
fermion bounds as stipulated in the primitive input in \cref{sec:inputs}.
\section{The primitive and nonlinear-step inputs}\label{sec:inputs}

We use the photon and fermion constructions of
\cite{MaxwellBlockManuscript,FermionBlockManuscript}, the extraction
and reconstruction maps of \cite{NormalizedContraction}, and the
nonlinear maps of \cite{NonlinearStep}.
We recall their hypotheses and the conclusions needed here.
The finite expressions used in the comparison are recalled in \cref{app:formulas}.

For an active step put \(e=e_k\), \(e_{k+1}=\sqrt L e_k\) and
\(0<\mu\le L^{-1}\). Both charges lie below the common ceiling of the
primitive estimates. The incoming and outgoing sides are admissible:
\(\ell/(LM)\) and \(\ell/(L^2M)\) are odd integers at least \(13\).
Assume \(0<\rho<1\), odd \(L\ge4^8\),
\(M=L^m\ge L^2\), \(M>4L\),
\(0<\varepsilon\le\bar\varepsilon=1/200\), and
\(1-\varepsilon<\alpha<1\). The condition \(\ell\ge Mr_k\)
is imposed at every active step.

The constants in the block estimates are
\[
 \begin{gathered}
 c_X=54,\qquad \kappa_0=54\log72,\qquad K_{\mathrm{pol}}=\exp(\kappa_0)/36,\qquad
 c_{\rm rb}=c_XK_{\mathrm{pol}},\qquad c_{\rm fl}=61c_{\rm rb},\\
 c_C=3L^{9/2},\qquad b=1,\qquad b_L=b/L,\qquad b_-=1-L^{-1},\\
 C_T=2\exp(3/2)
       [1+b_-^{-1}(54L\exp(6)+12\exp(1)+1)],\\
 c_S=10^5c_X^5K_{\mathrm{pol}}^5\exp(6\kappa),\qquad
 \kappa_w=\kappa+\kappa_0+3.
 \end{gathered}
\]
The constants \(C_w,\gamma_w,\alpha_0,C_B\) and \(c_\gamma\) are those
of the primitive families; \(c_\gamma,C_w\ge1\).
Here \(c_\gamma\) bounds the exponential row sums on the photon
coordinate and unit-bond index sets. The determinant constant is
\[
 C_d=16c_XC_w\max\{\tfrac\pi2 b\exp(1)(1+2\exp(1)),
                         2+58\exp(1)+b\exp(1)\}.
\]
For related tori \(C_w\) always means the common constant after
projection, \(64K_{\mathrm{pol}}C_w^{(1)}\).

\paragraph{Primitive hypotheses.}
Use the full hypotheses and conclusions of
\cite[Theorems A--B and Proposition 5.1]{MaxwellBlockManuscript}
and \cite[Theorem A and Propositions 4.1--4.3]{FermionBlockManuscript}.
In particular the parameter families are defined on
\(|s_{\rm pc}|\le M^{\alpha_0}\), with common analytic kernel and
short H\"older bounds, exact full-strength values, component
restriction, the row-vanishing properties stated for each family,
and local read bonds.
Write $s=(s',s'')$, where $s'=(s'_M,s'_{LM})$ consists of
the photon reconstruction and square-root parameters, of kinds
$\pi_M,\pi_{LM}$ in \cite[(37)]{MaxwellBlockManuscript}.
The fermion parameters $s''$ are indexed by the partition
$\pi_k$ into $M$-cubes, as in
\cite[\S4.1]{FermionBlockManuscript}.
For a photon site $c$, put
$K_c=\{c':|c'-c|_\infty\le2L\}$; the row of $H_k$ at $(c,g)$
vanishes unless every $M$-cube meeting $K_c$ is active.
Their lattice, gauge, transpose and reality identities are retained.
At zero use the diagonal families and the full chamber covariance.
The local resolvent bound then retains its factor \((1+y)^{-1}\).
The top fermion estimates are taken with polymer decay
\(\kappa_w+\kappa_0\); the projected estimates have decay \(\kappa_w\).

The finite inverse and walk conditions of the fermion construction
and the photon weakening condition
\[
 0<\alpha_0\le\tfrac14,\qquad
 2\alpha_0 n_P\log M\le\gamma_PM,\qquad n_P=402
\]
are assumed. Exact photon intertwining under odd covers is
\cite[(46)]{MaxwellBlockManuscript}. The comparison below proves
intertwining of the exact fermion formulas directly.

\paragraph{Extraction, reconstruction and determinant estimates.}
Use \cite[Theorems A--B and Appendix B]{NormalizedContraction},
including the fine inverses at both signs of the total connection
required by the reconstruction theorem.
Use the weighted integration convention of
\cite[Theorem B and Proposition 5.1]{FermionBlockManuscript}
and the determinant estimates of
\cite[Propositions 6.1--6.2]{FermionBlockManuscript}.
The extraction maps are linear, real on real activities, and satisfy
\[
 |m(\mathscr E)|\le c_{\rm fl}e_k^{1/2}\|\mathscr E\|_k^x,\qquad
 |\varepsilon_{\rm v}(\mathscr E)|\le c_{\rm fl}\|\mathscr E\|_k^x .
\]
The normalized linear term and determinant term satisfy
\[
 \|\mathcal L(R\mathscr E)^{\rm nat}\|_{k+1}^x
      \le c_{\rm fl}L^{-1/4+2\varepsilon}\|\mathscr E\|_k^x,
 \qquad
 \|\mathcal L\delta E_k^{\rm det}\|_{k+1}^x
      \le e_{k+1}^{1/4-2\varepsilon}.
\]
The scaled domains obey
\[
 \mathscr A_L\in L^{-3/4}\widetilde R_k\subset\widetilde R_k^-,
 \qquad \mathscr C_L\in L^{-3/4}\mathfrak X_k.
\]
The complete real activity spaces, all these maps and their norms
commute with restriction to the external slot \(\mathscr C=0\);
this restriction does not increase the norm.

\paragraph{Numerical conditions for the nonlinear step.}
The hypotheses of \cite[Theorem A]{NonlinearStep} include
\[
 \begin{gathered}
 L^{1/8-3\bar\varepsilon/2}\ge2400,\qquad e^\varepsilon\le\tfrac14,\qquad
 160L^{1+3\varepsilon/2}Me^\varepsilon\le1,\\
 9ML^\varepsilon e^{2\varepsilon}\le1,\qquad
 90c_XLMe^\varepsilon\le1,\qquad
 \kappa\ge7\kappa_0+7,\qquad
 L(\kappa-6\kappa_0-4)-1\ge\kappa,
 \end{gathered}
\]
and
\[
 C_w\ge4\exp(2\gamma_w)M^{27\alpha_0},\qquad
 M^{\alpha_0}-1\ge\exp((2L+1)^3\kappa+\log(6c_X)).
\]
For \(v=-\log e\), \(p_{0,k}=v^{p_0}\) and \(p_k=v^p\), impose
\[
 \begin{gathered}
 p_{0,k}\le e^{-\varepsilon/4},\qquad
 12\exp(-p_{0,k}^2/2)\le e^7,\qquad 2e^{5\varepsilon/2}\le1,\qquad
 c_{\rm fl}e^\varepsilon\le1,\qquad 3Me^\varepsilon\le1,\\
 2c_CC_w^2c_\gamma^2e^{\varepsilon/4}\le1,\qquad
 16c_Xc_\gamma C_wC_T^2M^3e^{\varepsilon/2}\le1,\qquad
 96M^3e^{1/4}\le1,\\
 32C_wC_Te^{3\varepsilon/2}\le1,\qquad
 88M^3e^{1/4-6\varepsilon}\le1,\qquad
 112C_w^2c_\gamma e^{2\varepsilon}\le3,\\
 c_Se^{1/4-5\varepsilon}\le1,\qquad c_SL^3e^\varepsilon\le1,\qquad
 8c_{\rm rb}L^3C_dM^3e^{2\varepsilon}\le1,\\
 8C_Bp_k\le e^{-3/4+3\varepsilon},\qquad
 8e^{3/4-7\varepsilon/2}\le1,\qquad L^{3/4}\ge8/3,\qquad
 v\ge9,\qquad 8M(v/4+4)e^\varepsilon\le1,\\
 e^{1/4}+e^\varepsilon<1,\qquad
 4e^{1/4+2\varepsilon}+2e^{3\varepsilon}\le1.
 \end{gathered}
\]
These conditions are imposed at each active scale. The additional
conditions used for the stopping arithmetic and the finite-sequence
contraction are stated with the physical scales.

\paragraph{The imported nonlinear conclusions.}
The displayed finite contour definitions use
\[
 r_1=\tfrac12e^{-3/4+9\varepsilon/2},\qquad
 r_2=e^{-3/4+7\varepsilon},\qquad r_4=e^{-1/4+3\varepsilon}.
\]
For a bond field $Z$ put
$|Z|_* =\max\{\|Z\|_\infty,\|\partial Z\|_\infty,
\|\delta_\alpha\partial Z\|_\infty\}$.
By \cite[Proposition 5.1]{NonlinearStep},
\(2\le r_2,r_4\le r_1\) and
\[
 |Z(s')|_*\le\tfrac12e^{-\varepsilon/2},\qquad
 A+tZ(s')\in\widetilde R_k,\qquad
 A+\mathscr C+tZ(s')\in\widetilde{\mathcal G}_k
       \quad(|t|\le r_1).
\]
The finite activities have the local ordinary, pair and fluctuation
supports and read bonds stated in
\cite[Proposition 5.3]{NonlinearStep}.
Their connected union expansion converges absolutely in the weighted
norm and its exponential is the original finite fluctuation integral,
by \cite[Proposition 6.1]{NonlinearStep}.

For every complex datum in \(\operatorname{Ball}_k^x(4)\), the resulting
correction satisfies
\[
 \|E_k^*(m,\mathscr E)\|_{k+1}^x\le e_k^{1/4-6\varepsilon}.
\]
It is norm analytic on affine complex lines in the interior of the
allowed data ball, continuous on their closed allowed sets, and real
for real data. For two data in \(\operatorname{Ball}_k^x(1)\),
\[
 \|E_k^*(m,\mathscr E)-E_k^*(\widetilde m,\widetilde{\mathscr E})\|_{k+1}^x
 \le e_k^{1/4-6\varepsilon}
 \left(\frac{|m-\widetilde m|}{w_k^{\rm m}}
       +\frac{\|\mathscr E-\widetilde{\mathscr E}\|_k^x}{w_k^{\rm E}}\right).
\]
At \(k=0\) the activity is fixed to zero and the second quotient is
absent. These are \cite[Proposition 8.1]{NonlinearStep}.
The full mass, vacuum, activity and external-scalar recursions reproduce
the exact density and its unnormalized photon-box marginal by
\cite[Theorem A]{NonlinearStep}.
In particular \(0\le\delta\varepsilon_k^0\le e_k^7\).
All these conclusions commute with restriction to \(\mathscr C=0\).


The exact external scalar update, with $\mathfrak P_0=0$, is
\begin{equation}\label{eq:rgnl-external-scalar-map}
 \mathfrak P_{k+1}(Y,\mathscr C')
  =\sum_{\bar X^L=LY}
       [\mathfrak P_k(X,\mathscr C'_L)+E_k^{\mathrm{det}}(X,\mathscr C'_L)].
\end{equation}
The full activity update is
\[
 \mathscr E^+=\mathcal L(R\mathscr E)^{\mathrm{nat}}
                +\mathcal L\delta E_k^{\mathrm{det}}+E_k^*(m,\mathscr E),
 \qquad
 \delta E_k^{\mathrm{det}}(X;A,\mathscr C)
       =E_k^d(X,A+\mathscr C)-E_k^d(X,\mathscr C).
\]
With no external argument the determinant difference is
$E_k^{\mathrm{det}}(X,A)=E_k^d(X,A)-E_k^d(X,0)$.
At the first step the incoming activity vanishes. There are
$192\ell_0^3$ incoming chamber pairs of measure $1/48$ and
$4\ell_1^3$ outgoing pairs. The fluctuation covariance is $48\delta$
on the chamber modes, and the local vertex is
\eqref{eq:rg-step-first-vertex}.
\section{Projection of the primitive systems}\label{sec:rg-cover-comparison}

We construct the primitive systems on one large torus and
project them to smaller tori through odd covering maps. Projection preserves the
estimates with a fixed loss in the common constant and
one loss in polymer decay. The comparison of local steps
then uses supported rows near a common lift, where the
signed image sum has only one surviving term.
This is the construction of \cite[\S\S15.1--15.3]{QED3Master}.

\subsection{Odd covers and projected kernels}
\label{subsec:rgcov-projection}

Fix $k\ge0$, $0<\mu\le1$, $0<e_k\le\bar e_0$, and admissible
sides $\ell_2=J_c\ell_1$, where $J_c$ is odd.
Write $T_i=\eta_k\Z^3/\ell_i\Z^3$ and
$p:T_2\to T_1$ for the cover. Let
$\operatorname{pr}_i$ denote the covering-lattice projections.
Pullback $p^*$ is inclusion of functions on the covering
index sets. They have character
$\chi(w)=1$ for photon fields, coordinates and parameters,
and $\chi(w)=(-1)^{w_1+w_2+w_3}$ for spinors and pair
fields. Oddness gives
$\chi(J_cw)=\chi(w)$, so a field on $T_1$ is a field
on $T_2$ with the same character.
The image is
\[
 \mathrm{Per}=\{f:f(t+\ell_1w)=\chi(w)f(t),\ w\in\Z^3\}.
\]
Translations move the photon site and preserve the chamber
label. In particular pullback commutes with $J,J^*$,
complex conjugation and constant spin matrices.

The index weight is $1$ on unit sites, unit bonds and
photon coordinate indices, $\eta_k^3$ on fine photon
sites and bonds, and $\eta_k^3/48$ on chamber sites.
For fields with a common index set and character,
\begin{equation}\label{eq:rgcov-pullback-pairing}
 \|p^*f\|_\infty=\|f\|_\infty,\qquad
 \langle p^*f',p^*f\rangle_{T_2}
                                      =J_c^3\langle f',f\rangle_{T_1}.
\end{equation}
Indeed a fundamental set for $T_2$ is the union of
$\mathrm{Fnd}_1+\ell_1w$, over representatives
$w\in(\Z/J_c\Z)^3$. The two characters cancel in
each paired summand, and the weight is the same on both
tori. At zero the chamber weight remains $1/48$.

Pullback maps the parameter polydisc, including its interior,
into the corresponding larger polydisc, since
$(p^*s)_{\mathrm{pc}}=s_{p\,\mathrm{pc}}$.
It also maps
$\widetilde{\mathcal G}_k(T_1)$ into
$\widetilde{\mathcal G}_k(T_2)$.
To check the latter assertion, write the domain decomposition
$A=A_0+A_1$. Pullback preserves reality of $A_0$ and
the sup norm of $A_1$. Each five-cube neighborhood has
a common nonwrapping lift of side $11M<\ell_1$.
The real local gauge representative of $A_0$ on its
image in $T_1$ pulls back on this lift to the
corresponding neighborhood in $T_2$, with its same
sup bound and gradient identity. This is precisely
the local real-gauge condition in the inverse domain.

Every exact finite fermion operator used below intertwines
with pullback. For the chamber Dirac operator this follows
term by term from its four neighbors and bond phases.
For $\mathcal Q_k=Q_kJ^*$ and its reverse it follows from
the same block sums and the same six paths on the cover;
$J^*$ averages the same $48$ labels.
The penalty and the finite matrices multiplying the
resolvent have the same finite block formulas.
Sums and products preserve intertwining.
If $O_2p^*=p^*O_1$ and both operators are invertible,
then
$O_2p^*O_1^{-1}=p^*$ and hence
$O_2^{-1}p^*=p^*O_1^{-1}$.
The domain preservation just proved makes the already
established inverse bounds applicable on both tori.
It follows, from the formulas in Section~\extref{4}{FermionBlockManuscript},
that the fine resolvents, $H_k,\bar H_k,D_k,M_k,\Gamma_k$
intertwine at the same connection:
\begin{equation}\label{eq:rgcov-exact-fermion-intertwining}
 O_2(p^*A)p^*=p^*O_1(A).
\end{equation}
At zero use the entire chamber Dirac stencil and the
inverse of $F_y$, on all chamber components.
The exact photon operators have the corresponding
identity \exteq{46}{MaxwellBlockManuscript}.

Let a primitive system on $T_2$ be given.
For a parameter family $O_2(s,A)$ put
$\widetilde O(s,A)=O_2(p^*s,p^*A)$.
For a polymer family put
\[
 \widetilde O(X,A)=\sum_{X'\in\mathcal P_k(T_2):pX'=X}
                                                    O_2(X',p^*A).
\]
This is applied to $M_k$ for $k\ge1$ and to $S_{k,y}$
at every level. A projected polymer is connected:
adjacent $M$-cubes have equal or adjacent images.
Moreover
\begin{equation}\label{eq:rgcov-projected-polymer-length}
                           d_M(pX')\le d_M(X').
\end{equation}
A projected net meets all image cubes. Projecting its chosen
finite segment representation preserves the lengths of the segments;
taking the infimum gives \eqref{eq:rgcov-projected-polymer-length}.
There are only finitely many polymers on either torus.

\begin{proposition}\label{prop:rgcov-kernel-projection}
If the system has translation covariance under
$\ell_1\Z^3$, its projection is well defined by
$p^*O_1=\widetilde Op^*$ and has kernel
\begin{equation}\label{eq:rgcov-signed-image-kernel}
 O_1(\cdot)(t,t')
  =\sum_{w\in(\Z/J_c\Z)^3}\chi(w)
                       \widetilde O(\cdot)(t,t'+\ell_1w).
\end{equation}
This formula uses the same index weights on both tori.
The signed summand depends only on the class of $w$.
Transposition of an intertwining operator also intertwines:
\begin{equation}\label{eq:rgcov-transpose-intertwining}
                          O_2^Tp^*=p^*O_1^T.
\end{equation}
\end{proposition}

\begin{proof}
The pulled-back parameters and connection are invariant
under $\ell_1$ translations. Parameter-family covariance
therefore says that $\widetilde O$ commutes with those
translations. For polymer families, translation permutes
the polymers $X'$ with $pX'=X$, so the same commutation
holds after their finite sum.
An operator commuting with these translations maps the
character eigenspace $\mathrm{Per}$ into itself.
Since $p^*$ is injective onto that space, this defines
exactly one projected operator.

Sum its action over the fundamental sets used in
\eqref{eq:rgcov-pullback-pairing} and insert
$f(t'+\ell_1w)=\chi(w)f(t')$.
The coefficient of $\nu(t')f(t')$ is
\eqref{eq:rgcov-signed-image-kernel}.
Replacing $w$ by $w+J_cv$ multiplies the character
by $\chi(v)$ and the large-torus kernel by the same
character, so the product is unchanged.
This proves the asserted kernel formula, including
the weights $\eta_k^3/48$ or $1/48$ on chamber columns.

Translations have transpose $T_v^T=T_{-v}$ for the
displayed bilinear pairings. Hence $O_2^T$ also commutes
with $\ell_1$ translations and maps $\mathrm{Per}$
into itself. Write $O_2^Tp^*f'=p^*g$.
For every source field $f$, intertwining and
\eqref{eq:rgcov-pullback-pairing} give
\[
 J_c^3\langle g,f\rangle_{T_1}
 =\langle O_2^Tp^*f',p^*f\rangle_{T_2}
 =J_c^3\langle f',O_1f\rangle_{T_1}.
\]
Nondegeneracy gives $g=O_1^Tf'$.
\end{proof}

\subsection{The uniform projection estimate}
\label{subsec:rgcov-projection-estimates}

The supports of primitive row and column indices have
sup diameter at most $L$. Assume
$\gamma_wM\ge1$. Then
\begin{equation}\label{eq:rgcov-exponential-image-sum}
 \sum_{w\in(\Z/J_c\Z)^3}
     \exp\big(-\gamma_wd_2(t,t'+\ell_1w)\big)
                     \le64\exp\big(-\gamma_wd_1(t,t')\big).
\end{equation}
To prove this, translate the column lift by an
$\ell_1$ period so that the distance of the two support
sets equals $d_1\le\ell_1/2$.
For each class of $w$, choose a representative $w'$
which realizes its distance on $T_2$. Distinct classes
have distinct representatives. Their distances are
at least $d_1$; for $|w'|_\infty=j\ge2$ they are at least
\[
 \ell_1j-d_1-2L
                         \ge d_1+(\ell_1-2L)(j-1).
\]
This uses the support diameters and $d_1\le\ell_1/2$.
Since $\ell_1\ge13LM$, we have
$\gamma_w(\ell_1-2L)\ge6$.
There are $27$ lattice vectors of norm at most $1$
and at most $26j^2$ of norm $j$.
The sum is therefore bounded by
\[
 \exp(-\gamma_wd_1)
                 \left(27+\sum_{j\ge2}26j^2\exp(-6(j-1))\right).
\]
The latter series starts below $0.269$ and its
successive ratios are at most
$(9/4)\exp(-6)<1/100$. Thus the bracket is below
$27.3<64$, proving \eqref{eq:rgcov-exponential-image-sum}.

There is a second, polymer summation estimate. For a
unit or chamber row index $t$, let $\square_t$ be its
$M$-cube on $T_2$. By
\eqref{eq:rgcov-projected-polymer-length} and counting,
\begin{equation}\label{eq:rgcov-polymer-image-sum}
 \sum_{\substack{pX'=X\\X'\supseteq\square_t}}
    \exp\big(-(\kappa_w+\kappa_0)d_M(X')\big)
                 \le K_{\mathrm{pol}}\exp\big(-\kappa_wd_M(X)\big).
\end{equation}
Indeed the factor with rate $\kappa_w$ is at most
the factor on the right; the remaining sum is bounded
by $\sum_{X'\supseteq\square_t}\exp(-\kappa_0d_M(X'))\le K_{\mathrm{pol}}$.
The additional top-system decay $\kappa_0$ pays for the
polymer image sum.

\begin{proposition}\label{prop:rgcov-projected-system}
Suppose the system on $T_2$ has the primitive properties
with parameters
\[
 (\alpha_0,\gamma_w,\kappa_w+\kappa_0;C_w^{(1)}),
 \qquad C_w^{(1)}\ge4\exp(2\gamma_w)M^{27\alpha_0}.
\]
Its projection has the corresponding properties on $T_1$
with
\begin{equation}\label{eq:rgcov-common-constant}
                          C_w=64K_{\mathrm{pol}}C_w^{(1)}
\end{equation}
and parameters $(\alpha_0,\gamma_w,\kappa_w;C_w)$.
In particular it satisfies the primitive hypotheses of
Theorem~\extref{A}{NonlinearStep}.
At zero the resolvent estimate retains $(1+y)^{-1}$.
\end{proposition}

\begin{proof}
We verify the identities, estimates and local structure
needed for the step.
At full strength the exact intertwining identities and
the primitive sum identities on $T_2$ give
\[
 H_{k,1}^{\mathrm{ph}}(1)=H_{k,1}^{\mathrm{ph}},\quad
 C_{k,1}^{1/2}(1)=C_{k,1}^{1/2},\quad
 H_{k,1}(1,A)=H_{k,1}(A),\quad
 \Gamma_{k,1}(1,A)=\Gamma_{k,1}(A).
\]
For the two polymer families, sum first over all $X$:
each $X'$ occurs once, under $X=pX'$.
Equation~\eqref{eq:rgcov-exact-fermion-intertwining}
then identifies the sum with the exact operator on $T_1$.
All projected entries are finite sums of original entries
composed with the linear pullback. Continuity on the
closed parameter polydisc, holomorphy in its interior and
in the connection for any fixed parameter, and Borel
dependence on the resolvent parameter follow from these
finite sums.

For the photon reconstruction, its first and H\"older
derivative kernels have the same image-sum formula:
the difference acts in the row variable, and short
lifted differences agree on the two tori.
Equation~\eqref{eq:rgcov-exponential-image-sum} bounds
each by $64C_w^{(1)}\exp(-\gamma_wd_1)$.
It gives the same bound for the square-root kernel
and for $\Gamma_k$.
The $H_k$ sup and short-pair H\"older bounds instead
follow by the isometric pullback in
\eqref{eq:rgcov-pullback-pairing}; their constant is
$C_w^{(1)}$. A short pair and its six paths lift
identically. An intracluster quotient is
$\eta_k^{-\alpha}(f(c,g')-f(c,g))$ on each torus.

For $M_k(X,A)$, only $X'$ containing $\square_t$
contribute to its row, by original polymer support.
Combine \eqref{eq:rgcov-signed-image-kernel},
\eqref{eq:rgcov-exponential-image-sum} and
\eqref{eq:rgcov-polymer-image-sum} to get
\[
 |M_{k,1}(X,A)(t,t')|_c
     \le C_w\exp(-\gamma_wd_1(t,t')-\kappa_wd_M(X)).
\]
For a positive-level fine resolvent, the pullback of
a unit-block indicator is the disjoint sum of its
$J_c^3$ translated large-torus block indicators.
The same two estimates bound its block operator
norm. At zero use its chamber kernel directly:
the original factor $(1+y)^{-1}$ is common to all
terms of both image sums and remains on the right.
All column weights are unchanged, so no further
factor $48$ appears.

For component locality, use shared-site linked components for
the photon parameter cubes and face-connected components for
the fermion partition $M$-cubes. An index whose support is
contained in its active union belongs to exactly one such
component. Fermion site supports are singletons, each in one
partition cube. For the photon kinds, the geometric argument
is as follows. For a bond of
length $a\in\{\eta_k,1\}$, cubes of radius $M$ meeting
its support either share an endpoint or would have
centers differing in that direction by $2M+a$.
The latter is impossible since the centers differ by
a multiple of $M$ and $0<a<M$.
For a coordinate index, its support has diameter at
most $L$. The centers of two cubes of radius $LM$
meeting it differ coordinatewise by less than $3LM$,
and are multiples of $LM$; choosing the coordinate
midpoint when the difference is $2LM$, and either
center when it is $LM$, exhibits a lattice site in
both cubes. Thus all cubes meeting the support are
linked. The same argument applies on either torus
by lifting the support and these cubes into a
nonwrapping neighborhood.

Pull back the active parameter cubes on $T_1$.
A component on $T_2$ maps into one component on $T_1$,
since the covering map preserves shared sites for photon
cubes and common faces for fermion partition cubes.
For a row in a small-torus component, the unique
large-torus component containing its lift maps into
that component. Original component restriction
confines all surviving columns there.
Changing the parameters and connection outside the
small component leaves these original rows unchanged;
the signed image sum therefore preserves their row
dependence. The identical argument applied to each
lifted column preserves column component restriction.
Each row-vanishing property stated for an original family passes
to its projection. In particular, if a cube meeting $K_c$ is
missing, its unique lift meeting the lifted $K_c$ is missing
too, so the $H_k$ row vanishes. For $k\ge1$, the
parameter-independent term $\mathcal B_{k,0}$ in $\Gamma_k$
is retained; it preserves every union of partition $M$-cubes
\cite[Proposition~4.3 and its proof]{FermionBlockManuscript}.
For a projected polymer family a nonzero term has
both original indices in $X'$, so both projected
indices lie in $X$. Equality of connections on bonds
meeting $X$ gives equality on bonds meeting every
such $X'$, preserving the polymer's local dependence.

Gauge covariance passes through pullback because
the local gradients, restriction to unit sites and
phases commute with it. Lattice covariance follows
by lifting a symmetry $(S,a)$ with one representative
$a\in LM\Z^3$ to the same affine map on both tori;
it permutes the image polymers. A change of representative
multiplies source and target spinors by the same sign,
which cancels in the operator relation.
Transposition passes by
\eqref{eq:rgcov-transpose-intertwining};
the constant $\mathsf C$ commutes with pullback.
Complex conjugation and $\varkappa$ also commute
with it. These facts prove the charge and reality
relations for the projected primitives.

At zero, the photon and fermion diagonal factors
are unchanged by projection. Cubes containing a
given bond, and cubes meeting $K_c$, have unique
corresponding lifts, so their parameter products
are the same on both tori. The fermion factor is
the same for all $48$ labels at $c$.
Their constants are bounded by
$4\exp(2\gamma_w)M^{27\alpha_0}\le C_w^{(1)}$.
The remaining zero-level properties are those
just proved for the square root, chamber covariance
and local resolvent. This completes the system
verification with the common constant
\eqref{eq:rgcov-common-constant}.
\end{proof}

\subsection{A common top-torus system}
\label{subsec:rgcov-top-system}

On a chosen top torus use the photon families
\exteq{37}{MaxwellBlockManuscript} and the local
fermion families \exteq{57}{FermionBlockManuscript}.
In the fermion walk construction take the polymer
decay rate $\kappa'=\kappa_w+\kappa_0$.
The formulas do not depend on this rate; it specifies
the parameter restrictions and the constants in their
estimates. We write $\gamma^F=\gamma$ for the fermion
decay rate in \exteq{31}{FermionBlockManuscript}. Define
\begin{gather}
 \alpha_0=1/54,\qquad
 \gamma_w=\min(\gamma^F,\gamma^P),\quad
                                \gamma^P=\gamma_P/2, \notag\\
 C^P=116C_P^{\mathrm{ph}}M^{402\alpha_0},\qquad
 C_w^{(1)}=\max\{C^F,C^P,4\exp(2\gamma_w)M^{27\alpha_0}\},
 \qquad C_w=64K_{\mathrm{pol}}C_w^{(1)}.             \label{eq:rgcov-top-constants}
\end{gather}
Here $C^F,\gamma^F$ are the constants of the fermion
construction at the indicated polymer rate.
The fixed $M$ satisfies both its fermion condition and
\exteq{38}{MaxwellBlockManuscript}; the charge satisfies
the inverse and walk conditions of those constructions.

These families form a system with
$(\alpha_0,\gamma_w,\kappa_w+\kappa_0;C_w^{(1)})$.
This follows by matching their proved properties.
The photon kernel multipliers are polynomials in their
parameter cubes; at full strength the partition of unity
and the products give the exact kernels. Their path
products confine rows and columns to their active
link-components. The photon derivative and square-root
bounds have constants $C^P$ and rate $\gamma^P$.
The fermion families are the convergent weakened walks
and the localized block expressions, with their proved
support, covariance, analyticity and H\"older estimates.
Replacing their rates $\gamma^F,\gamma^P$ by the smaller
$\gamma_w$ and their constants by the larger $C_w^{(1)}$
preserves every inequality. The row and column
restriction identities contain no decay constant and
are unchanged.

At level zero the two diagonal families are used.
The photon derivative bound has constant
$4\exp(2\gamma_w)M^{27\alpha_0}$, and the fermion
diagonal factor is bounded by
$M^{8\alpha_0}$, below that constant.
There is no primitive short pair with $0<d_X<1$ at $\eta_0=1$.
The covariance and local resolvent are chamber kernels
with measure $1/48$, and the resolvent bound includes
$(1+y)^{-1}$. A unit block has total chamber measure
$1$, since
$(\eta_k^3/48)(48\eta_k^{-3})=1$.
Consequently the positive-level block estimates and
zero-level chamber estimates use precisely the common
constants \eqref{eq:rgcov-top-constants}.
Projecting this system by \cref{prop:rgcov-projected-system}
gives a system with these primitive bounds and the same $C_w$ on every
odd smaller torus. The top system itself has these
weaker bounds as well. None of these constants depends
on the level, side, mass, charge or harmonic background.

\section{Agreement of local steps}\label{sec:local-comparison}

\subsection{Common lifts and local algebras}
\label{subsec:rgcov-common-lifts}

For the rest of the comparison assume
\begin{equation}\label{eq:rgcov-comparison-side}
                      \ell_1>(31L+16)M.
\end{equation}
Put $d_U=2L+1$, $c_l=2L$.
For a finite nonempty union $\widetilde X$ of covering
$M$-cubes, let $\operatorname{box}(\widetilde X)$ be
its smallest rectangular union of cubes and let
$\operatorname{diam}(\widetilde X)$ be that box's
largest side. Define
\[
 N(\widetilde X)=\operatorname{box}(\widetilde X)^{\sim6d_U}.
\]
We call $\widetilde X$ admissible
when $\operatorname{diam}(\widetilde X)
\le\ell_1-14d_UM$; a polymer is flat when it is the
image of a connected admissible $\widetilde X$.
These are the sets on which the two torus algebras
will be identified.

\begin{proposition}\label{prop:rgcov-common-lifts}
For admissible $\widetilde X$, projection is injective
on $N(\widetilde X)$ and on every supported index
contained there. An index of either torus supported
in the image of a subset of $N(\widetilde X)$ has
exactly one lift contained in that subset.
Pairs have their shortest lifts, with the same
$d_X$ and six paths.
Local incidences of total support diameter below
$2d_UM$ are preserved in both directions.

There is a norm-preserving algebra isomorphism
$\iota_{\widetilde X}$ between the two subalgebras
generated by ordinary fields over $N(\widetilde X)$,
pair fields with both endpoints there, and their
unit fluctuation and auxiliary fields there.
It identifies covering generator indices.
At zero it includes the chamber $W$ fields and
no auxiliary $\Psi$.
\end{proposition}

\begin{proof}
Every side of $N(\widetilde X)$ is at most
$\ell_1-2d_UM$. Since faces contain no fine lattice
sites, any two of its sites satisfy
\[
 |u-u'|_\infty\le\ell_1-2d_UM-\eta_k.
\]
For $w\ne0$ this implies
\begin{equation}\label{eq:rgcov-lift-separation}
              |u-u'-\ell_iw|_\infty\ge2d_UM+\eta_k.
\end{equation}
It proves disjointness of $N(\widetilde X)$ from
its nonzero $\ell_1$ translates, injectivity, and
uniqueness of any representative with distance
below $2d_UM$.

To lift an index supported in the projected subset,
translate one of its lifts so that one support point
lies in that subset. Every remaining support point
is congruent to a point of the subset and is within
$2LM<2d_UM$ of the first. Equation~
\eqref{eq:rgcov-lift-separation} forces the congruence
translation for each point to be zero.
Thus the whole index lies in the subset; injectivity
gives uniqueness. In particular the shortest pair
has its endpoints there, distance at most $M$, and
its same covering paths.

For an incidence relation, lift a tuple realizing
the projected relation and translate its first index
to the prescribed lift. Each other lifted index
differs from the prescribed one by an $\ell_i$
period. Its support is at distance below $2d_UM$
from the first one. Equation~
\eqref{eq:rgcov-lift-separation} again forces the
period to vanish. This proves the reverse incidence
implication; the forward one is projection.
The relations used below meet this range condition:
nearest neighbors and chamber neighbors, short
pairs and their paths, block and cube membership,
common faces, linked parameter cubes, cube meeting,
and coordinate support. The largest diameter is
$4LM<2d_UM$, for two linked photon parameter cubes.
They are all invariant under $LM\Z^3$ translations.

No two covering generator indices in $N(\widetilde X)$
have the same period class on either torus.
Relative to any fundamental generator system their
images are therefore independent generators, up
to their antiperiodic signs.
Identifying the same covering indices extends
uniquely to a free Grassmann algebra isomorphism.
It preserves bar indices and degrees and sends
basis monomials to basis monomials, up to the
same coefficient signs. All generator weights and
index measures are common, so it preserves norms.
The maps for nested admissible sets restrict to
one another.
\end{proof}

A local activity on $X_i=\operatorname{pr}_i\widetilde X$
belongs to this identified algebra. Ordinary sites
lift to $\widetilde X$; an admissible pair lifts
to $\widetilde X^{\sim1}$. For its first kind the
two sites and paths lie in $\widetilde X$.
For its second kind the pair is $(x,x^{(3)})$
with an endpoint over $\widetilde X$.
Each assertion follows from the incidence comparison.
Two local families $F_1,F_2$ agree on flat polymers
when
\begin{equation}\label{eq:rgcov-local-agreement}
 \iota_{\widetilde X}F_1(X_1;A_1,\mathscr C_1)
                             =F_2(X_2;A_2,\mathscr C_2)
\end{equation}
whenever the two admissible connection pairs have
equal lifts on $N(\widetilde X)$.
It suffices to check $A_2=p^*A_1$,
$\mathscr C_2=p^*\mathscr C_1$:
for other pairs with equal lifts, locality of
$F_2$ identifies them with these pullbacks on all
bonds meeting $\widetilde X$.
Thus once agreement is proved, equality holds
already under equality on those read bonds.

\begin{theorem}\label{thm:rgcov-step-comparison}
Suppose the two sides at levels $k$ and $k+1$ are
admissible and satisfy
$\ell_1/L>(31L+16)M$. Use the projected and top
systems constructed above, with the hypotheses of
Theorem~\extref{A}{NonlinearStep} on both tori.
Let the real mass $m$ be common with $|m|\le w_k^{\mathrm m}$,
and let the incoming real activities satisfy
$\|\mathscr E_i\|_k^x\le w_k^{\mathrm E}$ and agree on flat polymers,
with $\mathscr E_i=0$ at zero. For the external scalar conclusion,
assume the incoming scalar families agree on flat polymers as well.
Then their extracted mass and vacuum contributions
are equal. Their outgoing activities agree on flat
polymers at level $k+1$.
The external scalar recursion
\eqref{eq:rgnl-external-scalar-map} preserves agreement
of scalar families on such polymers.
All these assertions hold for the activities without
an external slot as well.
\end{theorem}

The proof is given in \cref{subsec:rgcov-step-comparison}, after
comparison of the local operations.

Every small polymer is flat. The net diameter bound
gives
$\operatorname{diam}(\widetilde X)
\le M(d_M(\widetilde X)+3)\le M(L+3)$ for a small
connected lift. By \eqref{eq:rgcov-comparison-side},
$M(L+3)+14d_UM<\ell_1$.
The nonwrapping lift proposition then identifies
small torus polymers containing a chosen cube
with small covering polymers containing its
chosen lift. This identification is the same
on both tori. Conversely, if an admissible
connected lift projects to a small polymer,
the small nonwrapping lift is its period translate:
start at one common cube and propagate equality
along adjacent cubes using the incidence result.
Hence smallness, tree length and volume coincide
for its two images.

Subpolymers and parameter cubes contained in an
admissible $\widetilde Y$ correspond bijectively
on the two tori. Common faces and links are
preserved, so incompatibility is preserved.
The definitions by cube meeting give
\[
 P_0(\operatorname{pr}_i\widetilde X)
                 =\operatorname{pr}_iP_0(\widetilde X),\qquad
 \operatorname{pr}_iY'(\widetilde X,\widetilde P)
          =Y'(\operatorname{pr}_i\widetilde X,\operatorname{pr}_i\widetilde P).
\]
When $\widetilde X$ and all cubes of $\widetilde P$
are contained in $\widetilde Y$, the resulting
$Y'$ is contained there too, since the $M$-cubes
partition sites. Reblocking satisfies
$\overline{\widetilde X}^{\,L}
\subset\widetilde Y^{\sim(L-1)}$ and commutes with
projection. These larger sets still lie in
$N(\widetilde Y)$, so their incidences have the
same unique lifts.

\subsection{Supported rows near a flat set}
\label{subsec:rgcov-supported-rows}

The preceding identification applies to the weakened
primitives because setting parameters to zero separates
their components. We first record the component argument.
Let $s$ vanish at every parameter cube not contained
in a union $Z$ of $M$-cubes. If $Z_1$ is a face-connected
component of $Z$ and
$\square(t)^{\sim2L}\subset Z_1$, then
\begin{equation}\label{eq:rgcov-component-row}
 \begin{split}
 O(s,A)(t,t')&=0 \quad\text{unless }t'\text{ is contained in }Z_1,\\
 O(s,A)(t,\cdot)&=O(s\mathbf1_{Z_1},A')(t,\cdot).
 \end{split}
\end{equation}
Here $O$ is a parameter primitive, and $A'$ is any
connection in the primitive's domain equal to $A$ on bonds meeting $Z_1$;
the photon primitives have no connection argument.
The set $\square(t)$ is the union of the $M$-cubes
meeting the support of $t$.

To prove \eqref{eq:rgcov-component-row}, choose a parameter
cube containing that support inside
$\square(t)^{\sim2L}$. For a chamber site or a unit site
it is its $M$-cube. For a fine bond it is the parameter
cube of radius $M$ about the $M$-cube center of one end:
both ends lie in it, and it is contained in
$\square(t)^{\sim1}$. For a photon coordinate index choose
the parameter cube of radius $LM$ about the $LM$-cube
center of one point of its support.
The support diameter is at most $L$ and $M\ge2$,
so this cube contains the entire support.
If $z_c$ is the $M$-cube center of that point and $z''$
is the $M$-cube center of any site in the parameter cube,
then
\[
 |z''-z_c|_\infty
 <(3L/2+1)M\le(2L+1)M.
\]
These centers differ by a multiple of $M$, hence the
cube is contained in $\square(t)^{\sim2L}$.

Let $\mathscr Y$ be all parameter cubes of the chosen kind
contained in $Z$, and let $\mathscr Y_1$ be the component
containing the cube just chosen, using shared-site links for
photon cubes and face adjacency for fermion partition cubes.
The $M$-cubes meeting
$\bigcup\mathscr Y_1$ form a face-connected polymer
contained in $Z$. It meets $\square(t)\subset Z_1$,
and therefore is contained in the component $Z_1$.
The component restriction property of the weakened
primitives, proved in the photon and fermion constructions,
now confines the row to $\bigcup\mathscr Y_1$.
Replacing $s$ by $s\mathbf1_{Z_1}$ leaves all its
parameters there unchanged; replacing $A$ by $A'$
leaves all its read bonds unchanged. This proves
\eqref{eq:rgcov-component-row}.

\begin{proposition}\label{prop:rgcov-supported-rows}
Let $\widetilde Y$ be admissible, $Y_i=\operatorname{pr}_i
\widetilde Y$, and $0\le j\le5d_U$. Let $s_1$ vanish
outside $Y_1$. Define $s_2$ by giving the corresponding
parameter cubes contained in $\widetilde Y$ the same
values and setting all other entries to zero.
Suppose $A_i\in\widetilde{\mathcal G}_k(T_i)$ have equal
lifts on the bonds contained in
$\widetilde Y^{\sim(j+d_U)}$.
For each parameter primitive and each covering row $t$
contained in $\widetilde Y^{\sim j}$, the two rows,
written with covering indices in $N(\widetilde Y)$,
coincide:
\begin{equation}\label{eq:rgcov-local-primitive-row}
                 O_1(s_1,A_1)(t,t')=O_2(s_2,A_2)(t,t').
\end{equation}
Both rows have columns contained in
$Y_i^{\sim(j+c_l)}$, and read the connection only on
bonds contained in $Y_i^{\sim(j+d_U)}$.

For a connected nonempty $\widetilde X\subset\widetilde Y$
and connections with equal lifts on bonds meeting
$\widetilde X$, the local polymer families also satisfy
\begin{equation}\label{eq:rgcov-local-polymer-kernel}
 O_1(\operatorname{pr}_1\widetilde X,A_1)(t,t')
       =O_2(\operatorname{pr}_2\widetilde X,A_2)(t,t')
\end{equation}
at rows and columns over $\widetilde X$.
Here $O=M_k$ for $k\ge1$, or $O=S_{k,y}$ for any $k$
and $y\ge0$.
\end{proposition}

\begin{proof}
Put $\widetilde B=\widetilde Y^{\sim(j+c_l)}$.
Its one-cube enlargement is
$\widetilde Y^{\sim(j+d_U)}\subset N(\widetilde Y)$.
For the row $t$, the set $\square(t)$ is one $M$-cube
or two face-adjacent $M$-cubes. This follows directly
for site and bond indices; a coordinate support lies
in one $L$-block or is the support of one bond.
Thus $\square(t)^{\sim c_l}$ is a connected box
contained in $\widetilde B$.
Let $\widetilde B_1$ be the component of $\widetilde B$
containing this box.

Its image is a component of
$\operatorname{pr}_1\widetilde B$.
On $T_2$ its image is a component both of
$\operatorname{pr}_2\widetilde B$ and of
$p^{-1}(\operatorname{pr}_1\widetilde B)$.
Indeed a cube adjacent to its image has a lift adjacent
to a cube of $\widetilde B_1$. This lift lies in
$\widetilde B^{\sim1}\subset N(\widetilde Y)$.
If it belongs to a period translate of $\widetilde B$,
separation \eqref{eq:rgcov-lift-separation} forces that
translate to be zero. It then lies in $\widetilde B$
and in the same component. This proves maximality
of the images as components.

Apply \eqref{eq:rgcov-component-row} on $T_1$.
It confines the row to $\operatorname{pr}_1\widetilde B_1$
and lets the connection be changed off the bonds meeting
that component. Such bonds lift into
$\widetilde B^{\sim1}$, proving the claimed read set.
On the parameter cubes contained in
$\operatorname{pr}_2\widetilde B_1$, the families
$p^*s_1$ and $s_2$ agree. To see this even at cubes
not contained in $\widetilde Y$, observe that each
such cube in $N(\widetilde Y)$ has a unique lift;
its image cannot also be the image of a cube in
$\widetilde Y$. Both entries are therefore zero.
The connections $p^*A_1$ and $A_2$ agree on bonds
meeting $\operatorname{pr}_2\widetilde B_1$.
Applying \eqref{eq:rgcov-component-row} twice on $T_2$
gives
\[
 O_2(p^*s_1,p^*A_1)(t,\cdot)
       =O_2(s_2,A_2)(t,\cdot),
\]
with columns confined to $\operatorname{pr}_2\widetilde B_1$.

In the signed image formula
\eqref{eq:rgcov-signed-image-kernel}, a nonzero summand
at $t'+\ell_1w$ would have a lift in $\widetilde B_1$.
Both this lift and $t'$ belong to $N(\widetilde Y)$;
their difference is an $\ell_1$ period.
Separation forces $w=0$ in $(\Z/J_c\Z)^3$.
That summand has character $1$, proving
\eqref{eq:rgcov-local-primitive-row}.

For the polymer families, every polymer $X'$ on $T_2$
with $pX'=\operatorname{pr}_1\widetilde X$ is
$\operatorname{pr}_2(\widetilde X+v)$ for one
$v\in\ell_1\Z^3/\ell_2\Z^3$.
Distinct classes give disjoint cube sets.
To verify the assertion for an arbitrary $X'$, lift
two adjacent cubes as cubes of $\widetilde X+v_1$
and $\widetilde X+v_2$. Their adjacency places a
translate of the second cube in
$\widetilde X^{\sim1}\subset N(\widetilde Y)$.
Separation forces $v_1=v_2$ modulo $\ell_2\Z^3$.
Connectedness gives one class for every cube of $X'$,
and projection onto the whole image gives equality
with that entire copy.

The projected kernel is consequently a double sum over
$v$ and the column image $w$. The local polymer support
requires both its row and column to be over
$\operatorname{pr}_2(\widetilde X+v)$.
The common row lift forces $v=0$; the common column
lift then forces $w=0$. The remaining term is
$O_2(\operatorname{pr}_2\widetilde X,p^*A_1)(t,t')$.
Locality on the bonds meeting this polymer permits
$p^*A_1$ to be replaced by $A_2$.

The barred mean field satisfies
$\bar H_k(s,A)=\mathsf C^{-1}H_k(s,-A)\mathsf C$,
so its comparison follows at the connections $-A_i$.
At zero the mean-field kernels are
$48\delta_{x,x'}\mathbf1_4$ times their parameter products;
the coordinate and chamber measures are the same
on the two tori, and the chamber label is preserved.
There is no $M_0$ family. The proof applies to
$\Gamma_0$ and $S_{0,y}$ on the full chamber indices.
\end{proof}

\subsection{Finite rows and the local operations}
\label{subsec:rgcov-finite-rows}

For finite operators, equality of rows follows directly
from their covering formulas. Suppose a row has all its
columns in $N(\widetilde Y)$, and its connection and
polymer membership data agree on the bonds and sites
which that row reads. The two formulas are then the
same finite sum on either torus. There are no repeated
column classes by \cref{prop:rgcov-common-lifts}.
Writing the row as
$\sum_{t'}r(t,t')f(t')$, its kernel is
\begin{equation}\label{eq:rgcov-row-kernel-weight}
                           O(t,t')=\nu(t')^{-1}r(t,t').
\end{equation}
Thus equality holds for the kernels as well, with the
same column weight.

We need this observation at connections $B_i$ having
equal lifts on bonds meeting a connected
$\widetilde U\subset\widetilde Y$.
The photon coordinate map $C$ reads no connection:
a noncentral face bond is its own coordinate, and
the central bond is the negative sum of the other
face bonds and the interior straight-path contributions
$n_c(b)/L$. All columns are supported in the one or
two $L$-blocks of the row bond's ends. Hence its rows
at bonds in $\widetilde U^{\sim(1+2d_U)}$ agree and
have range at most $2L-1$.

The covariant difference row at an admissible pair is
\[
 d_X(x,x')^{-\alpha}
               \big(U_B(x,x')f(x')-f(x)\big).
\]
Its only columns are $x,x'$, and it reads exactly
the six pair paths. These paths are contained in
$\widetilde U$ for the first kind, or are the bond
meeting $\widetilde U$ for the second kind.
The range is $M$; inside a chamber cluster the
transport is $1$ and no bond is read.
These rows therefore agree, also at $-B_i$.

For $k\ge1$ the reconstruction rows $T^{(1)},T^{(2)}$
and their barred versions are the two-piece formulas
of the reconstruction estimate. At $y\in\widetilde U$
they have columns over $\Delta_y$ and at $x^{(3)}$,
or at the pair $(x,x^{(3)})$ across its boundary.
They read the bonds in $\Delta_y$ and the bonds
$(c,c_3)$ starting there. All meet $\widetilde U$.
Their membership indicators are the same because
\[
 u\in\widetilde U+\ell_i\Z^3
       \quad\Longleftrightarrow\quad u\in\widetilde U
             \qquad(u\in\widetilde U^{\sim1}),
\]
again by separation. Every site involved is at
distance at most $(1+\eta_k)/2<1$ from $y$.
This proves equality of both pieces, including
the boundary term.

The finite rows of $Q_kJ^*$ and $JQ_k^T(-B)$ have
their columns and read bonds in a single $\Delta_y$.
The row of $B_{k,r}$, where $r\ge0$ is the resolvent
parameter, has its columns in one $L$-block and reads
the paths from its center to its sites.
These paths lie in that block. Their kernels are
\begin{align*}
 (Q_kJ^*)(y,x)&=\mathbf1[c\in\Delta_y]\varphi_B(y,c)\mathbf1_4,\\
 (JQ_k^T(-B))(x,y)&=\mathbf1[c\in\Delta_y]\varphi_B(y,c)^{-1}\mathbf1_4,\\
 B_{k,r}
   &=(b_k+\mathrm i\gamma_5r)^{-1}I\\
   &\quad+\big((b_k+b_L+\mathrm i\gamma_5r)^{-1}
                    -(b_k+\mathrm i\gamma_5r)^{-1}\big)P_B .
\end{align*}
The projection $P_B$ has kernel $L^{-3}$ times the two
block transports when its indices lie in the same
$L$-block, and zero otherwise. The first two rows
have range below $1$, the third range at most $L-1$.
At rows in $\widetilde U$ their columns and read
bonds stay in its $M$-cube. They agree on both tori.

At zero, $F_0$ has only its diagonal, the four chamber
neighbors, and its block-penalty columns. The hopping
part of its kernel is
\[
 48\Big(
  \bigl(\mu+\tfrac{\rho}{2}\sum_{i=0}^3w_i\bigr)\delta_{x,x'}\mathbf1_4
       -\sum_{j=0}^3\delta_{x',x^{(j)}}M_j(x)U_B(x,x^{(j)})
       \Big).
\]
The factor $48$ is the inverse column weight.
Only the fourth neighbor reads a bond meeting
$\widetilde U$; the penalty reads paths inside the
$L$-block of $c$. All columns lie over
$\widetilde U^{\sim1}$ and its range is at most $L-1$.
Thus this entire row agrees at zero.

These equalities also identify the mass form
with the covering sum:
\[
 \langle\bar\psi,\psi\rangle_{\operatorname{pr}_i\widetilde\square}
 =(\eta_k^3/48)\sum_{c\in\widetilde\square,g,\beta}
                          \bar\psi_\beta(c,g)\psi_\beta(c,g).
\]
At zero the two halves of the first-step vertex have
the common coefficient $1/96$, the same two boundary
transports, and the same generators over the cube
and at its outward pairs. Hence the first vertex
also agrees under $\iota_{\widetilde Y}$.

The remaining local operations can be compared in
this common free algebra. Equality of elements gives
equality of their constant coefficients and of their
antisymmetric quadratic kernels: expand in ordered
covering monomials and use the common measure
$\eta_k^3/48$. For a small polymer, the adjustment
$\Theta_{X,B}$ uses averages over shortest chains
of $M$-cubes in $X$. Common faces, path membership
and transports are identical on its lift, as are
its volume and the coefficients $d_X^\alpha$.
Therefore
\begin{equation}\label{eq:rgcov-adjustment-intertwining}
 \iota_{\widetilde X}\Theta_{X_1,B_1}
                      =\Theta_{X_2,B_2}\iota_{\widetilde X}.
\end{equation}

If $\widetilde U$ is admissible at level $k+1$,
then $L\widetilde U$ is admissible at level $k$.
Its box diameter is $L\operatorname{diam}(\widetilde U)$,
and $N(L\widetilde U)\subset L N(\widetilde U)$.
The scaling rules on the common generators are
$\psi(x)\mapsto L^{-1}\psi(x/L)$ and
$\chi(x,x')\mapsto L^{-1-\alpha}\chi(x/L,x'/L)$.
They preserve the chamber label and the period signs.
Thus
\begin{equation}\label{eq:rgcov-scaling-intertwining}
 S_L\iota_{L\widetilde U}=\iota_{\widetilde U}S_L.
\end{equation}

Finally, the integrations agree on any local element.
For $W$, expand in monomials. A monomial moment is
zero unless its modes are paired; otherwise it is
its ordering sign times $\nu^{-\#\mathrm{pairs}}$.
The mode indices and the weight $\nu$ agree.
At zero this weight is $1/48$, so the factor is
$48^{\#\mathrm{pairs}}$; other modes integrate to $1$.
This proves
\begin{equation}\label{eq:rgcov-W-integration}
 \iota_{\widetilde Y}\int\dd\mu_I(W)\,F
                    =\int\dd\mu_I(W)\,\iota_{\widetilde Y}F.
\end{equation}
The restricted photon law is a product of the same
normalized Gaussian on $[-p_{0,k},p_{0,k}]$ in each
coordinate. Consequently its marginal on common
coordinates is identical. With $d_V=2-2L^{-3}$,
its vacuum contribution on either torus is
\begin{equation}\label{eq:rgcov-box-vacuum}
 \delta\varepsilon_k^0
 =-d_V\log\left((2\pi)^{-1/2}
             \int_{-p_{0,k}}^{p_{0,k}}\exp(-z^2/2)\dd z\right).
\end{equation}
This number depends on $e_k$ and $L$, and is independent
of the side.

\subsection{Extraction and determinant comparison}
\label{subsec:rgcov-determinants}

Suppose $\mathscr E_i\in K_k^x(T_i)$ agree in the sense
of \eqref{eq:rgcov-local-agreement}. For a small
covering polymer $\widetilde X$, their values at $(0,0)$
have the same constant and quadratic coefficients.
Their local extractions
\eqref{eq:rg-local-extractions} are therefore the same:
their integrated kernels and their common volume are
the same. Fixing a covering cube identifies all small
polymers containing its image on either torus. The sums
\eqref{eq:rg-counterterm-maps} consequently give
\begin{equation}\label{eq:rgcov-counterterms}
 m(\mathscr E_1)=m(\mathscr E_2),\qquad
 \varepsilon_{\mathrm v}(\mathscr E_1)
                         =\varepsilon_{\mathrm v}(\mathscr E_2).
\end{equation}
Subtracting the common constant and mass forms preserves
local agreement of $R\mathscr E_i$.
On a small polymer, its quadratic part also agrees;
\eqref{eq:rgcov-adjustment-intertwining} proves agreement
of $(R\mathscr E_i)^{\mathrm{nat}}$.
On other polymers neither of these operations changes
the relevant element.

The determinant terms agree on every flat polymer,
not just on small ones. Let $X_i=\operatorname{pr}_i
\widetilde X$ and suppose $B_i$ have equal lifts on
bonds meeting $\widetilde X$.
For $k\ge1$, set
\[
 O_1=B_{k,r}\gamma_5B_{k,r},\qquad
 O_2=Q_kJ^*S_{k,r}(X_i,B_i)JQ_k^T(-B_i).
\]
The trace in the determinant integral is
\begin{equation}\label{eq:rgcov-local-determinant-trace}
 \operatorname{Tr}(O_1O_2)
       =\sum_{z_1,z_2\in X_i\cap\Z^3}
                      \tr\big(O_1(z_2,z_1)O_2(z_1,z_2)\big).
\end{equation}
Here and below $\operatorname{Tr}$ is the trace on the whole finite
function space, while $\tr$ is the spin trace.
Indeed
\begin{align*}
 O_2(z_1,z_2)
   &=\int_{\Delta_{z_1}}\dd x\int_{\Delta_{z_2}}\dd x'\,
       \varphi_{B_i}(z_1,c)S_{k,r}(X_i,B_i)(x,x')
                            \varphi_{B_i}(z_2,c')^{-1}.
\end{align*}
Local polymer support confines $x,x'$ to $X_i$,
and $\Delta_z$ lies in the $M$-cube of $z$.
Thus both unit sites in this trace lie in $X_i$.
The rows of $B_{k,r}$ stay in their own $L$-block,
so the intermediate sum in $O_1(z_2,z_1)$ also stays
in $X_i$. All sums may now be formed on the lift
$\widetilde X$; the two antiperiodic kernel signs
cancel in each trace summand.
The finite-row comparison and
\eqref{eq:rgcov-local-polymer-kernel} make every factor
the same on both tori.

At zero the required traces are
\begin{align}
 \operatorname{Tr}(F_0(B_i)S_{0,r}(X_i,B_i))
   &=\int_{X_i}\dd x\int_{X_i}\dd x'\,
          \tr\big(F_0(B_i)(x,x')S_{0,r}(X_i,B_i)(x',x)\big),
                                      \label{eq:rgcov-zero-trace}\\
 \operatorname{Tr}(\gamma_5S_{0,r}(X_i,B_i))
   &=\int_{X_i}\dd x\,\tr\big(\gamma_5S_{0,r}(X_i,B_i)(x,x)\big).
                                      \label{eq:rgcov-zero-trace-gamma}
\end{align}
Support again restricts all arguments to $X_i$.
Both are the same covering sums with measure $1/48$
and all $48$ chamber labels. The local $F_0$ and
$S_{0,r}$ kernels agree, so the traces agree for
every $r\ge0$.

These pointwise identities can be integrated in $r$.
The required domination follows
from the primitive estimates. Since $P_B^2=P_B$,
\[
 B_{k,r}\gamma_5B_{k,r}
  =\gamma_5\big((b_k+\mathrm i\gamma_5r)^{-2}(I-P_B)
                +(b_k+b_L+\mathrm i\gamma_5r)^{-2}P_B\big).
\]
In the component matrix norm,
$|(a+\mathrm i\gamma_5r)^{-1}|_c
\le\sqrt2(a^2+r^2)^{-1/2}$ for $a>0$.
Hence the first trace is bounded by a finite
$B_i,X_i$ dependent constant times
$(b_k^2+r^2)^{-1}$.
The other factor is uniformly bounded by the local
resolvent estimate and the two finite block transports.
This majorant is integrable because $b_k>0$.
At zero, $|S_{0,r}(X_i,B_i)(x,x')|_c
\le C_w(1+r)^{-1}$.
The first trace in \eqref{eq:rgcov-zero-trace} is
integrated against $\dd r/r$ on $[1,\infty)$,
and \eqref{eq:rgcov-zero-trace-gamma} against $\dd r$
on $[0,1]$; both are integrable.
The kernels are Borel in $r$, so these are
ordinary absolutely convergent scalar integrals.

The two unsymmetrized determinant primitives agree.
Applying the argument at $-B_i,\bar B_i,-\bar B_i$
identifies all four terms in its symmetrization.
It gives
\begin{equation}\label{eq:rgcov-determinant-agreement}
              E_k^{d,(1)}(X_1,B_1)=E_k^{d,(2)}(X_2,B_2).
\end{equation}
Subtracting the values at $\mathscr C_i$ or at zero
proves the corresponding equality for
$\delta E_k^{\mathrm{det}}$ and $E_k^{\mathrm{det}}$.

\subsection{Comparison of the local fluctuation sums}
\label{subsec:rgcov-fluctuation-comparison}

Let $m$ be real with $|m|\le w_k^{\mathrm m}$, let
$\|\mathscr E_i\|_k^x\le w_k^{\mathrm E}$, and assume
$\mathscr E_i=0$ at zero and local agreement as above.
Fix $A_1\in\widetilde R_k^-(T_1)$ and
$\mathscr C_1\in\mathfrak X_k(T_1)$, and put
$(A_2,\mathscr C_2)=p^*(A_1,\mathscr C_1)$.
All small-charge and primitive hypotheses of the
nonlinear step are imposed on both tori.
In particular the contour radii satisfy
$1<r_2\le r_1$, $r_4>1$, and the domain inclusions
of Proposition~\extref{5.1}{NonlinearStep} hold.
We compare the terms of that step on a connected
nonempty $\widetilde U\subset\widetilde Y$.

Assign common coordinate values
$z_\Upsilon\in[-p_{0,k},p_{0,k}]$ to the photon basis
elements contained in $\widetilde U$, and set all
other coordinate values to zero on each torus.
Assign common parameter values of modulus at most
$M^{\alpha_0}$ to the parameter cubes contained in
$\widetilde U$, and set all other parameters to zero.
These data will be denoted by $\zeta_i,s_i$.
They belong to the respective box and polydisc
because indices on the lift are injective.

The fine fluctuation fields have the same lifts on
$\widetilde U^{\sim(1+d_U)}$:
\begin{equation}\label{eq:rgcov-fine-fluctuation}
                         Z_1(s'_1)(b)=Z_2(s'_2)(b).
\end{equation}
To check the successive supports in their composition,
write it in kernels,
\[
 Z_i(s'_i)(b)
  =\sum_{b',\Upsilon,\Upsilon'}
    H_k^{\mathrm{ph},(i)}(s'_i)(b,b')
          C_i(b',\Upsilon)
            C_k^{1/2,(i)}(s'_i)(\Upsilon,\Upsilon')\zeta_i(\Upsilon').
\]
At such a row $b$, \cref{prop:rgcov-supported-rows}
with $j=1+d_U$ confines $b'$ to
$\widetilde U^{\sim2d_U}$ and identifies the row kernels.
The coordinate row has columns in the $M$-cubes of
the two ends of $b'$, hence in
$\widetilde U^{\sim(1+2d_U)}$; those finite rows agree.
The square-root row, with $j=1+2d_U$, confines
$\Upsilon'$ to $\widetilde U^{\sim3d_U}$.
All these sets lie in $N(\widetilde U)$, and the
last coordinate values agree there.
The two compositions are therefore the same finite
covering sum. At zero the first row is diagonal
and the sum over $b'$ has one term.

Put $B_i=A_i+\mathscr C_i$ and
$B_i(t)=B_i+tZ_i(s'_i)$.
For $|t|\le r_1$ these fields lie in
$\widetilde{\mathcal G}_k(T_i)$, and
$A_i+tZ_i(s'_i)\in\widetilde R_k(T_i)$.
They have equal lifts on the bonds just compared.
The fermion mean-field rows at sites at distance
at most $\eta_k$ from $\widetilde U$ consequently agree,
as do the $\Gamma_k$ rows at unit sites in
$\widetilde U$ and the $M_k(X_i,B_i(t))$ kernels
for lifted subpolymers $\widetilde X\subset\widetilde U$.
Here the parameters vanish outside $U_i$.
The row-vanishing and component argument in the proof
of Proposition~\extref{5.3}{NonlinearStep} confines
their columns to $U_i$ itself: a nonzero mean-field
row has its active parameter component attached to
that set; the $\Gamma_k$ row starting in this component
has the same confined columns.
For $k\ge1$, the parameter-independent term $\mathcal B_{k,0}$
in $\Gamma_k$ also preserves $U_i$, a union of partition
$M$-cubes.
Thus all these fields are polynomials in the common
generators over $\widetilde U$.

For example, at a relevant chamber site,
\[
 (H_k^{(i)}(s_i'',B_i(t))\Psi)_\beta(x)
   =\sum_{y\in\widetilde U\cap\Z^3,\beta'}
         H_k^{(i)}(s_i'',B_i(t))(x,y)_{\beta\beta'}
                                      \Psi(y,\beta',0)
\]
for $k\ge1$, with the analogous barred sum.
Their coefficients agree.
At every pair admissible for $\widetilde X$, applying
the common finite difference row preserves this equality.
The differences between these fields at $t=1$ and $t=0$
are precisely the ordinary and pair components of
$J_k^\#\Psi$ used in $G_2,G_3$.
Both values are in the disk because $r_1>1$.
Moreover locality of $R\mathscr E_i$ gives equality
of its values at the pairs
$(A_i+tZ_i,\mathscr C_i)$ on each subpolymer.
These pairs need only agree on its read bonds,
by the criterion following
\eqref{eq:rgcov-local-agreement}.

It follows that the three first-localization contour
integrands agree under $\iota_{\widetilde U}$.
For $G_1$, this is the common local $M_k$ bilinear.
For $G_2$, addition of the common odd field $tJ_k^\#\Psi$
to each generator commutes with an algebra homomorphism;
apply it to the equal local activity.
For $G_3$, the same substitution applies to the equal
mass form, and $m'=m+m(\mathscr E_i)$ is common by
\eqref{eq:rgcov-counterterms}.
Their contour radii and scalar contour measures
are the same. Integration thus gives
\begin{equation}\label{eq:rgcov-first-G}
              \iota_{\widetilde U}G_a^{(1)}(s_1,X_1)
                                      =G_a^{(2)}(s_2,X_2),
                    \qquad a=1,2,3.
\end{equation}
At zero $G_2=G_3=0$ and no auxiliary field occurs.
$G_1$ is the first-step vertex; its two factors
$\exp(\pm\mathrm i qZ_i(c,c_3))-1$ agree by
\eqref{eq:rgcov-fine-fluctuation}, with its common
coefficient $1/96$ and the common outward pair.
This proves \eqref{eq:rgcov-first-G} at that level too.

The finite localization commutes with this equality.
Indeed the pairs $(X,P)$ with $Y'(X,P)=U_i$ correspond
bijectively to lifted pairs $(\widetilde X,\widetilde P)$
inside $\widetilde U$ having
$Y'(\widetilde X,\widetilde P)=\widetilde U$.
If a torus pair has that union, its polymer and all
its parameter cubes are contained in $U_i$, so the
unique lifts give the inverse correspondence.
The identities for $P_0,Y'$ proved above show that
it preserves the parameter prescriptions.
Its summand is therefore
\begin{equation}\label{eq:rgcov-localization-term}
 \int_{[0,1]^{\widetilde P}}\dd u\,
       \prod_{\mathrm{pc}\in\widetilde P}\partial_{u_{\mathrm{pc}}}
                       G_a^{(i)}(s_i[u],\operatorname{pr}_i\widetilde X),\qquad a=1,2,3,
\end{equation}
where $s_i[u]$ is $1$ on $P_0(\widetilde X)$, $u$ on
$\widetilde P$, and zero elsewhere.
Identity \eqref{eq:rgcov-first-G} holds throughout
the parameter polydisc. Since $M^{\alpha_0}>1$,
the integration cube is in its interior, and
coefficient differentiation and integration preserve
the equality. This proves equality of
$\operatorname{Loc}G_a(U_i)$.

The reconstruction substitution replaces each $\Psi(y)$
by the same two-piece row $T^{(1)}\psi+T^{(2)}\chi$.
Both rows, including their boundary membership indicators,
have already been compared. Their resulting ordinary
generators are over $\widetilde U$, and their boundary
pairs are admissible for it. Equality on generators
proves equality of the substitution homomorphisms.
We obtain
\begin{equation}\label{eq:rgcov-first-localization}
              \iota_{\widetilde U}E_{\le3}^{(1)}(U_1)
                                      =E_{\le3}^{(2)}(U_2).
\end{equation}
At zero no reconstruction substitution is made.

For the second localization, the common fermion field is
\[
 \mathrm{Wc}=H_k(s'',B)\Gamma_k(s'',B)W,\qquad
 \overline{\mathrm{Wc}}=\bar H_k(s'',B)\bar W .
\]
For $k\ge1$ both are sums over the common unit sites
of $\widetilde U$. For $k=0$ they are
\[
 \mathrm{Wc}(x)=d(c)\frac1{48}
            \sum_{x'\text{ over }\widetilde U}
                                      \Gamma_0(s'',B)(x,x')W(x'),
 \qquad
 \overline{\mathrm{Wc}}(x)=d(c)\bar W(x),
\]
where the common diagonal factor $d(c)$ vanishes
outside $U_i$. These formulas preserve all chamber
labels and the weight $1/48$.
The kernel equalities proved above, followed by the
common pair difference rows, identify $W_k^\#$ at
every index used by a subpolymer.

The first-localized activity on a subpolymer reads
only photon coordinates supported there. Hence
\eqref{eq:rgcov-first-localization}, proved with coordinates
supported on that subpolymer, still holds when coordinates
are supported on the larger $\widetilde U$.
Adding the common $W_k^\#$ to its generators proves
equality of $G_{\le3}^{\,W}$.
The same addition to $R\mathscr E_i$ and to the mass
form proves equality of the contour integrands of
$G_4,G_5$. The second localization uses only $\pi_k$;
the correspondence and derivative argument
\eqref{eq:rgcov-localization-term} apply with this
one kind. Consequently
\begin{equation}\label{eq:rgcov-finite-F}
                       \iota_{\widetilde U}F_1(U_1)=F_2(U_2).
\end{equation}

We now pass to the local connected logarithm.
Its finite Mayer activity is
\[
 K_i(X_i)=\sum_{\{\widetilde X_a\}:\,\bigcup_a\widetilde X_a=\widetilde X}
                \prod_a\big(\exp(F_i(\operatorname{pr}_i
                                                    \widetilde X_a))-1\big),
\]
where the $\widetilde X_a$ are distinct connected
subpolymers of $\widetilde X$.
The products are even. Their index sets and factors
are the same on both tori by
\eqref{eq:rgcov-finite-F}; coordinate locality allows
their evaluation with common coordinates on
$\widetilde X$.
For the integrated activity $H_i(X_i)$,
\eqref{eq:rgcov-W-integration} identifies the Wick
integrals. The other photon coordinates integrate
to $1$, leaving the same product of restricted normal
laws over the coordinates supported in
$\widetilde X$. Hence $\iota_{\widetilde X}H_1(X_1)=H_2(X_2)$.

The logarithm has the expansion
\[
 E_k^{\#,(i)}(X_i)=
   \sum_{n\ge1}\frac1{n!}
    \sum_{\bigcup_{a=1}^n\widetilde X_a=\widetilde X}
       u^T(\operatorname{pr}_i\widetilde X_1,\ldots,
                    \operatorname{pr}_i\widetilde X_n)
                      \prod_{a=1}^n H_i(\operatorname{pr}_i\widetilde X_a).
\]
Projection preserves incompatibility inside this lift,
so the graph determining $u^T$ is the same.
Every summand agrees under the common algebra map.
Both series converge absolutely in the weighted norm
by Proposition~\extref{6.1}{NonlinearStep}.
The map $\iota_{\widetilde X}$ preserves that norm,
so equality passes to their sums:
\begin{equation}\label{eq:rgcov-logarithm-agreement}
       \iota_{\widetilde X}E_k^{\#,(1)}(X_1;A_1,\mathscr C_1)
                      =E_k^{\#,(2)}(X_2;A_2,\mathscr C_2).
\end{equation}

\subsection{Agreement of the step after reblocking}
\label{subsec:rgcov-step-comparison}



\begin{proof}[Proof of \Cref{thm:rgcov-step-comparison}]
The scalar conclusions are
\eqref{eq:rgcov-counterterms} and
\eqref{eq:rgcov-box-vacuum}.
For the outgoing activity it suffices to compare a
level $k+1$ connection pair with its pullback.
Let $\widetilde U$ be connected and admissible there,
put $Y_i=\operatorname{pr}_i'\widetilde U$,
and set $\widetilde Y=L\widetilde U$.
The latter is a connected union of $LM$-cubes,
hence also of $M$-cubes, admissible at level $k$.
Its neighborhood is contained in $L N(\widetilde U)$.

If $(\mathscr A',\mathscr C')$ satisfies the outgoing-domain hypotheses at level
$k+1$, its scaled pair lies in
$\widetilde R_k^-\times\mathfrak X_k$ by
\eqref{eq:rg-scaled-domain-inclusion}.
Scaling and pullback commute on covering bonds:
both give $L^{-1/2}B'(b/L)$.
Denote the scaled pairs by $(A_i,\mathscr C_i)$.
They satisfy exactly the hypotheses used in
\eqref{eq:rgcov-logarithm-agreement}.

The polymers with $\bar X^L=LY_i$ are precisely
$\operatorname{pr}_i\widetilde X$, where
$\widetilde X$ is connected in $\widetilde Y$ and
its $LM$-closure is $\widetilde Y$.
Indeed $X\subset\bar X^L$ gives its unique lift
inside $\widetilde Y$. Closure commutes with projection,
and lies in $\widetilde Y^{\sim(L-1)}$.
Projection is injective on that enlargement,
so equality of its projected closure with
$\operatorname{pr}_i\widetilde Y$ forces equality
of the lifted closure with $\widetilde Y$.
This also proves the converse and shows that the
reblocking index set is common.

In that index set the outgoing activity is
\begin{align}
 \mathscr E_i^+(Y_i;\mathscr A_i',\mathscr C_i')
  =\sum_{\substack{\widetilde X\subset\widetilde Y\text{ connected}\\
                         \overline{\widetilde X}^{\,L}=\widetilde Y}}
 S_L\Big[&
     (R\mathscr E_i)^{\mathrm{nat}}
                  (\operatorname{pr}_i\widetilde X;A_i,\mathscr C_i)
       +\delta E_k^{\mathrm{det},(i)}
                  (\operatorname{pr}_i\widetilde X;A_i,\mathscr C_i)
       \nonumber\\
       &+E_k^{\#,(i)}
                  (\operatorname{pr}_i\widetilde X;A_i,\mathscr C_i)\Big].
                                      \label{eq:rgcov-reblocked-comparison}
\end{align}
Each of the three entries agrees under
$\iota_{\widetilde Y}$, by extraction and adjustment,
\eqref{eq:rgcov-determinant-agreement}, and
\eqref{eq:rgcov-logarithm-agreement}, respectively.
The entries belong to its ordinary and pair subalgebra.
Equation~\eqref{eq:rgcov-scaling-intertwining} identifies
their scaled values, so the sums agree under
$\iota_{\widetilde U}$.
Locality of the outgoing family, already proved in
the step theorem, then extends the comparison
from pullback pairs to pairs with equal local lifts.

The external scalar recursion has the same reblocking
index set. Its incoming scalar term agrees by hypothesis,
and its local determinant term agrees by
\eqref{eq:rgcov-determinant-agreement} at
$(\mathscr C_1,\mathscr C_2)$ and at $(0,0)$.
Its two sums therefore agree.

The argument without an external slot uses precisely
these local operations with $B_i=A_i$.
The determinant difference is then
$E_k^{\mathrm{det}}(X_i,A_i)=E_k^d(X_i,A_i)-E_k^d(X_i,0)$.
Extraction, addition of odd fields, finite localization
and integration do not require that slot. Their
same domain and convergence assertions were established
for this case in the nonlinear-step construction.
Thus the three bracketed terms in
\eqref{eq:rgcov-reblocked-comparison} agree in that
case too, and scaling gives the conclusion.
Restriction of a family with a slot to
$\mathscr C=0$ also commutes with every operation
in the displayed formula.

At $k=0$ the incoming activity and its extracted
terms are zero. The comparison is for the local
determinant difference and the first correction,
with chamber generators, traces and Wick moments
of weight $1/48$, as explicitly included above.
\end{proof}


\section{Physical scales and stopping levels}\label{sec:rg-flow}

The logarithmic charge scale determines where the finite-volume
iteration stops. We express the stopping rule in terms of depth below
the physical mass scale. This removes the cutoff from the terminal
physical quantities and gives the side margins needed for comparison
at every preceding step.
The stopping-scale calculation follows \cite[\S16.1]{QED3Master}.

\subsection{Physical scales and stopping levels}
\label{subsec:rgfl-scales}

Fix $\bar m\in(0,1]$ and put
$h=L^{\lfloor\log_L(1/\bar m)\rfloor}$.
For a cutoff $N\ge0$, define
\begin{equation}\label{eq:rgfl-physical-scales}
 \eta=L^{-N},\qquad K=N+\log_Lh,\qquad
 a_j=L^{j-N},\qquad e_j=e\sqrt{a_j},\qquad
 \mu_j=\bar m a_j.
\end{equation}
Thus $a_j$ is the physical length of one unit at
level $j$, and $a_K=h$.
At zero the charge is written $q=e\sqrt\eta$;
the fixed charge ceilings $\bar e_0,\bar e_1$ retain their
earlier meanings. A formula with variable index $j$
uses $e_j$ for the running charge, including $q$
when $j=0$.
For $0\le j\le K$ set
\[
 v_j=-\log e_j,\qquad
 r_j=\min\{L^i:i\in\N\cup\{0\},\ L^i\ge v_j^r\},
 \qquad p_j=v_j^p,\quad p_{0,j}=v_j^{p_0}.
\]
The integers $r,p,p_0$ and the constants of the
preceding estimates are fixed.
Here $r\ge2$, $L$ is odd with $L\ge13$, and $M$
is a power of $L$ with $M\ge L^2$.
We assume the hypotheses of the primitive and
step estimates, and use a ceiling
$\bar e_1\le\min(\bar e_0,1/3)$ such that, for all
$v\ge-\log\bar e_1$,
\begin{equation}\label{eq:rgfl-scale-smallness}
 \begin{gathered}
 v\ge8,\qquad
 \left(1+\frac{3\log L}{2v}\right)^r<L,\qquad
 v^r>31L+16,\\
 (c_{\mathrm{fl}}+2)\exp(-\varepsilon v)
               \le\frac1{16}(1-L^{-\varepsilon/2}).
 \end{gathered}
\end{equation}
In addition
$0<\varepsilon\le1/200$,
$c_{\mathrm{fl}}\ge1$,
$\bar e_0^\varepsilon\le1/4$,
$c_{\mathrm{fl}}L^{-3/8+6/200}\le1/16$, and
$L^{-5/8}\le1/4$.
The construction below applies whenever
$0<e\sqrt h\le\bar e_1$.

The physical top side and reference side are
\begin{equation}\label{eq:rgfl-reference-sides}
 R_{\mathrm{top}}=hMr_K,\qquad
 R^*=R_{\mathrm{top}}/L=hL^{T^*(e)},\qquad
 T^*(e)=\log_L(Mr_K)-1.
\end{equation}
They are independent of $N$, since
$v_K=-\log(e\sqrt h)$.
The admissible physical sides for the flow are
\[
 \mathcal T(e)=\{hL^T:T\in\N\cup\{0\},\ hL^T\le R_{\mathrm{top}}\}.
\]
For $R\in\mathcal T(e)$ put
$\ell_j(R)=R/a_j$,
$T_j(R)=\eta_j\Z^3/\ell_j(R)\Z^3$ and
$\mathfrak n_j=\log_L(\ell_j/(Mr_j))$.
The allowed unit sides are
$\mathrm{Cl}=\{\ell:\ell/(LM)\text{ is an odd integer at least }13\}$.
Their occurrence at every required scale is part
of the following arithmetic argument.

\begin{proposition}\label{prop:rgfl-scales}
For each $R\in\mathcal T(e)$ there is a finite integer
$N_{\min}(R,e)\ge0$ such that, when $N\ge N_{\min}(R,e)$,
the first level with $\mathfrak n_j=0$ exists and
satisfies $1\le k_*(R)\le K$.
Its physical quantities $a_*,e_*,\mu_*,\ell_*,r_*$
are independent of $N$ in this range.
For $j\le k_*(R)$,
\begin{equation}\label{eq:rgfl-scale-ranges}
 \begin{gathered}
 \ell_j\in\mathrm{Cl},\qquad
 \ell_j\ge Mr_j,\qquad r_j\ge r_*>31L+16,\\
 v_j\ge v_*\ge v_K,\qquad
 0<e_j\le e_*\le e_K\le\bar e_1,\qquad
 0<\mu_j\le\mu_*\le\bar m h\le1 .
 \end{gathered}
\end{equation}
For $j<k_*(R)$ one has $\ell_j\ge LMr_j$,
$\ell_{j+1}=\ell_j/L\in\mathrm{Cl}$,
$e_{j+1}=\sqrt L e_j\le\bar e_0$ and
$0<\mu_j=\mu_{j+1}/L\le1/L$.
At the stop,
\begin{equation}\label{eq:rgfl-terminal-scales}
 \ell_*=Mr_*,\qquad
 Mv_*^r\le\ell_*<LMv_*^r,\qquad
 \mu_*^{-1}<L\ell_*,\qquad e_*^2\ell_*=e^2R.
\end{equation}
If $R\le R'$ in $\mathcal T(e)$, then
$N_{\min}(R',e)\le N_{\min}(R,e)$ and
$k_*(R)\le k_*(R')$ for a common cutoff above both stated thresholds.
\end{proposition}

\begin{proof}
Use the depth $s=K-j$. Its quantities are
\[
 \widehat v_s=v_K+\tfrac{s}{2}\log L,\qquad
 \widehat r_s=\min\{L^i:i\ge0,\ L^i\ge\widehat v_s^r\},\qquad
 \widehat n_s(R)=T+s-\log_L(M\widehat r_s).
\]
They do not depend on $N$. At depth $s=K-j$ they
equal $v_j,r_j,\mathfrak n_j$, while
$a_j=hL^{-s}$ and $\ell_j=L^{T+s}$.
Since $\widehat v_s\ge v_K$, the scalar condition gives
\[
 \widehat v_{s+1}^r
   =\widehat v_s^r\left(1+\frac{\log L}{2\widehat v_s}\right)^r
                    \le L\widehat v_s^r\le L\widehat r_s.
\]
Thus
$\widehat r_{s+1}\in\{\widehat r_s,L\widehat r_s\}$
and $\widehat n_{s+1}-\widehat n_s\in\{0,1\}$.
Also $\widehat n_0\le0$ because $R\le R_{\mathrm{top}}$.
Minimality of the rounded power gives
$\widehat r_s<L\widehat v_s^r$ and hence
\[
 \widehat n_s>T+s-1-\log_LM-r\log_L(v_K+\tfrac{s}{2}\log L)
                                                    \longrightarrow\infty .
\]
The integer sequence therefore has a largest zero.
Denote its index by $s_*(R,e)$.
It is nonpositive up to that index and at least $1$
afterwards.

The condition $\ell_0\ge LMr_0$ is
$\widehat n_K\ge1$, equivalently $K>s_*$.
Its exact threshold is
\begin{equation}\label{eq:rgfl-mesh-threshold}
 N_{\min}(R,e)=\max\{0,s_*(R,e)+1-\log_Lh\},
             \qquad k_*(R)=K-s_*(R,e).
\end{equation}
This gives $1\le k_*\le K$ and
$a_*=hL^{-s_*}$, proving cutoff independence of
the terminal quantities.
Increasing $T$ increases every $\widehat n_s$
by the same integer; thus it decreases $s_*$,
decreases $N_{\min}$ and increases $k_*$.

For $j\le k_*$, the depth satisfies $K-j\ge s_*$.
Monotonicity gives the asserted ranges of
$\mathfrak n_j,r_j,v_j,e_j$, and the running mass
increases with $j$.
Moreover
$\bar m h=L^{\lfloor\log_L(1/\bar m)\rfloor-\log_L(1/\bar m)}
\in(1/L,1]$.
Since $\mu_*\ell_*=\bar mR\ge\bar m h>1/L$,
we obtain $\mu_*^{-1}<L\ell_*$.
The rounded-power inequality at $v_*$ gives
the two bounds for $\ell_*$ in
\eqref{eq:rgfl-terminal-scales}; its last identity
uses $e_*^2=e^2a_*$.

Finally $r_*\ge v_*^r>31L+16>L$ implies
$r_*\ge L^2$, as it is a power of $L$.
At every required level,
$\ell_j/(LM)=r_jL^{\mathfrak n_j-1}$
is a positive power of the odd integer $L\ge13$.
This proves $\ell_j\in\mathrm{Cl}$.
Before the stop $\mathfrak n_j\ge1$, giving the
stronger side bound. The mass and charge bounds
for the step follow at $j+1$.
\end{proof}

On $R^*$ define the reference mesh threshold
\begin{equation}\label{eq:rgfl-reference-threshold}
                    N_{\mathrm{ref}}(e)=N_{\min}(R^*(e),e).
\end{equation}
The stop on this torus has depth $s_*\in\{1,2\}$.
Indeed, for $s=1,2,3$ the stronger inequality in
\eqref{eq:rgfl-scale-smallness} gives
$\widehat r_s\in\{r_K,Lr_K\}$.
On $R^*$,
$\widehat n_s=s-1-\log_L(\widehat r_s/r_K)$.
Its values at $s=0,1,2,3$ are respectively
$-1$, an element of $\{-1,0\}$, an element of
$\{0,1\}$, and an element of $\{1,2\}$.
The zero-step-or-one-step rule forces the last zero
to be $1$ or $2$. On $R_{\mathrm{top}}$ the sequence
is larger by one, so its last zero is $0$ or $1$.
Consequently
\begin{equation}\label{eq:rgfl-reference-stop}
 \begin{gathered}
 k_*(R^*)\in\{K-1,K-2\},\qquad
 a_*\in\{h/L,h/L^2\},\qquad r_*\in\{r_K,Lr_K\},\\
 k_{\mathrm{top}}\in\{K,K-1\},\qquad
 L^{-3}<\mu_*\le L^{-1}\quad\text{on }R^* .
 \end{gathered}
\end{equation}
All these physical terminal values are independent
of $N\ge N_{\mathrm{ref}}(e)$.

The scale proposition verifies every side, mass
and charge hypothesis of the step theorem
on $R_{\mathrm{top}}$ and on the smaller tori it covers with odd degree.
At each level use the top system of
\cref{subsec:rgcov-top-system}; on a smaller
torus use its projection from
\cref{prop:rgcov-projected-system}.
These are the actual mean-field, covariance,
local resolvent and photon families constructed
earlier, with the same common estimates.
The ratio $R_{\mathrm{top}}/R$ is a power of the
odd integer $L$, and $k_*(R)\le k_{\mathrm{top}}$.
Thus the required top families exist whenever
a smaller-volume step is used.
Furthermore
$\ell_{j+1}(R)\ge Mr_{j+1}>(31L+16)M$;
the next-side hypothesis of the cover comparison
holds at every such step.

\section{The terminal conditions and the common pair}\label{sec:common-tuning}

We give the continuation and tuning arguments of
\cite[\S\S16.2--16.4]{QED3Master}.

\subsection{Continuation within the mass range}
\label{subsec:rgfl-continuation}

Write the mass and activity bounds as
\[
 w_j^{\mathrm m}=e_j^{3/4-8\varepsilon},\qquad
 w_j^{\mathrm E}=e_j^{1/4-7\varepsilon}.
\]
For a real pair $(m_{\mathrm{in}},\varepsilon_{\mathrm{in}}^{\mathrm v})$,
the initial vacuum is $\varepsilon_{\mathrm{in}}^{\mathrm v}$,
the initial member of the mass sequence is $m_{\mathrm{in}}$,
and $E_0=\mathscr E_0=\mathfrak P_0=0$.
The variable-index notation $m_j$ includes this initial member.
The recursion is
\begin{align}
 m_{j+1}&=L(m_j+m(E_j)),&
 \varepsilon_{j+1}^{\mathrm v}
       &=L^3(\varepsilon_j^{\mathrm v}
                 +\varepsilon_{\mathrm v}(E_j)+\delta\varepsilon_j^0),
                                      \label{eq:rgfl-coupling-recursion}\\
 E_{j+1}&=\mathcal L(RE_j)^{\mathrm{nat}}
                          +\mathcal LE_j^{\mathrm{det}}+E_j^*(m_j,E_j),
                                      \label{eq:rgfl-activity-recursion}\\
 \mathscr E_{j+1}
       &=\mathcal L(R\mathscr E_j)^{\mathrm{nat}}
               +\mathcal L\delta E_j^{\mathrm{det}}
                                      +E_j^*(m_j,\mathscr E_j).
                                      \label{eq:rgfl-slot-recursion}
\end{align}
The scalar family $\mathfrak P_j$ follows
\eqref{eq:rgnl-external-scalar-map}.
In the last line the correction has the external slot;
in \eqref{eq:rgfl-activity-recursion} it has none.
These are the corresponding maps proved earlier.
The trajectory exists up to a level if every preceding
$(m_j,\mathscr E_j)$ lies in $\operatorname{Ball}_j^x(4)$.
It is controlled when it exists to $k_*(R)$ and
$|m_j|\le w_j^{\mathrm m}$ throughout, including the stop.

The normalized estimates used below are
\begin{equation}\label{eq:rgfl-window-estimates}
 \begin{gathered}
 L^{-1}\frac{w_{j+1}^{\mathrm m}}{w_j^{\mathrm m}}
                  =L^{-5/8-4\varepsilon}\le\tfrac14,\qquad
 c_{\mathrm{fl}}e_j^{1/2}
          \frac{w_j^{\mathrm E}}{w_j^{\mathrm m}}
                           =c_{\mathrm{fl}}e_j^\varepsilon,\\
 c_{\mathrm{fl}}L^{-1/4+2\varepsilon}
          \frac{w_j^{\mathrm E}}{w_{j+1}^{\mathrm E}}
                     =c_{\mathrm{fl}}L^{-3/8+11\varepsilon/2}
                                                        \le\tfrac1{16},\\
 \frac{e_{j+1}^{1/4-2\varepsilon}}{w_{j+1}^{\mathrm E}}
                                      \le e_{j+1}^\varepsilon,\qquad
 \frac{e_j^{1/4-6\varepsilon}}{w_{j+1}^{\mathrm E}}
             =L^{-1/8+3\varepsilon}e_{j+1}^\varepsilon
                                      \le e_{j+1}^\varepsilon .
 \end{gathered}
\end{equation}
Their first and third bounds use the fixed inequalities
stated above; their last two use $e_{j+1}<1$ and
$\varepsilon\le1/200$.
The charge condition also gives
\begin{equation}\label{eq:rgfl-summable-charge}
 \sum_{j=0}^{k_*}(c_{\mathrm{fl}}+2)e_j^\varepsilon
             \le(c_{\mathrm{fl}}+2)e_*^\varepsilon
                          (1-L^{-\varepsilon/2})^{-1}
                                                        \le\tfrac1{16}.
\end{equation}
In particular
$c_{\mathrm{fl}}e_j^\varepsilon\le1/16$,
$2e_j^\varepsilon\le1/16$, and
$e_j\le e_j^\varepsilon\le1/48$.

\begin{proposition}\label{prop:rgfl-continuation}
If the recursively produced masses satisfy
$|m_i|\le w_i^{\mathrm m}$ for $i<j\le k_*(R)$,
then the trajectory is defined to $j$ and
\begin{equation}\label{eq:rgfl-activity-bound}
 \|E_j\|_j\le\|\mathscr E_j\|_j^x
       \le\frac{32}{15}e_j^\varepsilon w_j^{\mathrm E}
       \le4e_j^{1/4-6\varepsilon}\le w_j^{\mathrm E}.
\end{equation}
The data satisfy
$E_j=\mathscr E_j|_{\mathscr C=0}$,
$\mathfrak P_j(X,0)=0$, and belong to their real
mass, vacuum and activity spaces.
\end{proposition}

\begin{proof}
At zero all activities and scalar families vanish.
Assume the assertion through $i<j$.
The mass condition and
$4e_i^\varepsilon\le1$ put the input in the unit ball,
so Theorem~\extref{A}{NonlinearStep} applies at $i$.
The scaled connection pair lies in
$\widetilde R_i^-\times\mathfrak X_i$, so every local
determinant and correction in the recursion is defined.
The scalar external sum is finite and uses
$\mathscr C_L\in\mathfrak X_i$.
The step theorem preserves the real spaces
and commutes with restriction to $\mathscr C=0$.
Consequently the coupling maps computed from
$E_i$ or $\mathscr E_i$ coincide.
The scalar recursion vanishes at zero because each
$E_i^{\mathrm{det}}(X,0)$ vanishes.

The activity bound of the current step, divided by
$w_{i+1}^{\mathrm E}$ and combined with
\eqref{eq:rgfl-window-estimates}, gives
\[
 \frac{\|\mathscr E_{i+1}\|_{i+1}^x}{w_{i+1}^{\mathrm E}}
      \le\frac1{16}
                   \frac{\|\mathscr E_i\|_i^x}{w_i^{\mathrm E}}
                                     +2e_{i+1}^\varepsilon.
\]
Starting from zero, the bound
$(32/15)e_i^\varepsilon$ is preserved since
$e_i\le e_{i+1}$ and $(32/15)/16+2=32/15$.
Restriction does not increase the activity norm.
This proves \eqref{eq:rgfl-activity-bound} and the induction.

At the first step the formulas use $192\ell_0^3$
incoming pairs with weight $1/48$ and
$4\ell_1^3$ outgoing pairs, as proved in the
exact density identity.
Its coupling maps are
$m_1=Lm_{\mathrm{in}}$ and
$\varepsilon_1^{\mathrm v}
 =L^3(\varepsilon_{\mathrm{in}}^{\mathrm v}+\delta\varepsilon_0^0)$;
its activity is the local determinant term plus
the first correction. The normalized estimate has
zero incoming-activity term and the same two
source bounds. Thus the induction includes that step.
\end{proof}

\subsection{The terminal conditions and the fixed point}
\label{subsec:rgfl-fixed-point}

\begin{theorem}\label{thm:rgfl-terminal-zero-pair}
On any $R\in\mathcal T(e)$ with $N\ge N_{\min}(R,e)$,
there is exactly one real controlled pair whose
trajectory has $m_{k_*}=0$ and
$\varepsilon_{k_*}^{\mathrm v}=0$.
Along its trajectory,
\begin{equation}\label{eq:rgfl-tuned-bounds}
 |m_j|\le\tfrac18w_j^{\mathrm m},\qquad
 \|E_j\|_j\le\|\mathscr E_j\|_j^x
                                  \le4e_j^{1/4-6\varepsilon},\qquad
 |\varepsilon_j^{\mathrm v}|
                          \le5c_{\mathrm{fl}}e_j^{1/4-6\varepsilon}.
\end{equation}
The vacuum at every level is
\begin{equation}\label{eq:rgfl-backward-vacuum}
 \varepsilon_j^{\mathrm v}
     =-\sum_{i=j}^{k_*-1}L^{-3(i-j)}
                      \big(\varepsilon_{\mathrm v}(E_i)
                                              +\delta\varepsilon_i^0\big).
\end{equation}
\end{theorem}

\begin{proof}
Take the finite product of the real mass coordinates
and real activity Banach spaces, with norm
\[
 \|\mathbf z\|_{\mathrm{fl}}
   =\max_{0\le j\le k_*}
       \max\left(\frac{|m_j|}{w_j^{\mathrm m}},
                              \frac{\|E_j\|_j}{w_j^{\mathrm E}}\right),
                  \qquad \mathbf z=(m_j,E_j)_{0\le j\le k_*}.
\]
The closed subset $\mathcal B$ defined by these
quotients being at most $1$ and $E_0=0$
is nonempty and complete. Define a self-map by
\begin{align}
 \widehat m_j&=L^{-1}m_{j+1}-m(E_j)\quad(j<k_*),
                 &\widehat m_{k_*}&=0, \nonumber\\
 \widehat E_0&=0,\qquad
 \widehat E_{j+1}
       =\mathcal L(RE_j)^{\mathrm{nat}}
                    +\mathcal LE_j^{\mathrm{det}}+E_j^*(m_j,E_j).
                                      \label{eq:rgfl-fixed-point-map}
\end{align}
All maps here are the actual maps without an external
slot; their inputs belong to the unit ball.
Extraction, adjustment followed by reblocking,
the determinant activity and the nonlinear correction
are defined on that ball and preserve the real space.

We verify that the map preserves $\mathcal B$.
The extraction estimate implies
$|m(E_j)|\le c_{\mathrm{fl}}e_j^{1/2}\|E_j\|_j$.
Equations~\eqref{eq:rgfl-window-estimates} and
\eqref{eq:rgfl-summable-charge} give
\[
 \frac{|\widehat m_j|}{w_j^{\mathrm m}}
                    \le\tfrac14+c_{\mathrm{fl}}e_j^\varepsilon
                    \le\tfrac5{16}.
\]
For the activity, the contraction, determinant and
correction estimates give
\[
 \frac{\|\widehat E_{j+1}\|_{j+1}}{w_{j+1}^{\mathrm E}}
                         \le\tfrac1{16}+2e_{j+1}^\varepsilon
                         \le\tfrac18 .
\]
The fixed end components are zero. Hence the map
takes $\mathcal B$ into itself.

For two sequences, let $d_j^{\mathrm m},d_j^{\mathrm E}$
be the corresponding normalized differences.
Linearity of extraction gives
\[
 \widehat d_j^{\mathrm m}
      \le\tfrac14d_{j+1}^{\mathrm m}
                            +c_{\mathrm{fl}}e_j^\varepsilon d_j^{\mathrm E}
      \le\tfrac5{16}\|\mathbf z-\widetilde{\mathbf z}\|_{\mathrm{fl}} .
\]
Adjustment followed by reblocking is linear in
the activity at fixed connection.
The local determinant term is independent of the
sequence. The correction's coupling estimate
\exteq{93}{NonlinearStep} therefore gives
\[
 \widehat d_{j+1}^{\mathrm E}
      \le\tfrac1{16}d_j^{\mathrm E}
                         +e_{j+1}^\varepsilon
                                      (d_j^{\mathrm m}+d_j^{\mathrm E})
      \le\tfrac18\|\mathbf z-\widetilde{\mathbf z}\|_{\mathrm{fl}} .
\]
The whole map contracts by $5/16$.
Banach's fixed point theorem
\cite[Theorem~8]{BrightMetric2023} gives a unique
fixed point in $\mathcal B$.
Its equations are precisely the mass and activity
recursions with $m_{k_*}=0$.

Let $m_{\mathrm{in}}$ be its initial mass coordinate
and choose
\[
 \varepsilon_{\mathrm{in}}^{\mathrm v}
      =-\sum_{i=0}^{k_*-1}L^{-3i}
                     \big(\varepsilon_{\mathrm v}(E_i)
                                           +\delta\varepsilon_i^0\big).
\]
This pair is real. Its full trajectory, including
the external slot, agrees with the fixed-point mass
and activity at every level: agreement at zero is
the initial condition; at the next level
\cref{prop:rgfl-continuation} makes the
full step defined, and the fixed-point equations
identify its mass and restricted activity.
The induction reaches $k_*$ and proves control.
Solving the scalar vacuum recursion backwards gives
\[
 \varepsilon_j^{\mathrm v}
 =L^{-3(k_*-j)}\varepsilon_{k_*}^{\mathrm v}
     -\sum_{i=j}^{k_*-1}L^{-3(i-j)}
                     \big(\varepsilon_{\mathrm v}(E_i)
                                           +\delta\varepsilon_i^0\big).
\]
At $j=0$ its comparison with the chosen input forces
$\varepsilon_{k_*}^{\mathrm v}=0$, giving
\eqref{eq:rgfl-backward-vacuum}.
Conversely any controlled pair with the two terminal
zeros has its mass/activity sequence in $\mathcal B$
by the continuation proposition. It is the same fixed
point. Its initial vacuum is then the same scalar
sum. This proves uniqueness of the pair.

For the sharper mass bound put
$B_j=|m_j|/w_j^{\mathrm m}$.
The backward mass equation and the unit activity
bound imply $B_j\le B_{j+1}/4+1/16$, while
$B_{k_*}=0$.
Summing gives $B_j\le(1/16)\sum_{s\ge0}4^{-s}
=1/12\le1/8$.
The activity bound in \eqref{eq:rgfl-tuned-bounds}
is the continuation estimate.

It remains to bound the vacuum sum.
Extraction gives
$|\varepsilon_{\mathrm v}(E_i)|
\le4c_{\mathrm{fl}}e_i^{1/4-6\varepsilon}$.
For $i\ge j$ the associated summand obeys
\[
 L^{-3(i-j)}e_i^{1/4-6\varepsilon}
       =e_j^{1/4-6\varepsilon}
                         L^{-(i-j)(23/8+3\varepsilon)}
       \le e_j^{1/4-6\varepsilon}16^{-(i-j)}.
\]
Thus its sum is at most
$(64/15)c_{\mathrm{fl}}e_j^{1/4-6\varepsilon}$.
The box normalization satisfies
$0\le\delta\varepsilon_i^0\le e_i^7$ by the
step estimate.
Its remaining sum is bounded by
\[
 \sum_{i=j}^{k_*-1}L^{-3(i-j)}e_i^7
        =e_j^6\sum_{i=j}^{k_*-1}e_i
        \le e_j^6 e_*(\sqrt L-1)^{-1}
        \le\tfrac1{48}e_j^{1/4-6\varepsilon}.
\]
Here $\sqrt L\ge2$ and $e_*\le1/48$.
Since $c_{\mathrm{fl}}\ge1$,
$64/15+1/48=343/80<5$, giving the claimed bound.
The sums are empty at the stop.

At zero $E_0=0$. The backward mass equation reads
$m_{\mathrm{in}}=L^{-1}m_1$, and the forward activity is
$\mathcal LE_0^{\mathrm{det}}+E_0^*(m_{\mathrm{in}},0)$.
The determinant and correction estimates used above
include their full chamber trace and the first
weighted fluctuation integral. The initial vacuum
sum contains $\delta\varepsilon_0^0$.
No additional coupling is selected by the fixed point.
\end{proof}

\subsection{One pair on all tori of the construction}
\label{subsec:rgfl-volume-comparison}

For $N\ge N_{\min}(R_{\mathrm{top}},e)$ define
\[
 P_N(e)=\big(m_{0,N}(e),\varepsilon_{0,N}^{\mathrm v}(e)\big)
\]
to be the pair of \cref{thm:rgfl-terminal-zero-pair}
on the prescribed side $R_{\mathrm{top}}(e)$.
The dependence on the fixed $\bar m$ and fixed
construction parameters is suppressed.
For lower cutoffs set $m_{0,N}(e)=0$.

\begin{theorem}\label{thm:rgfl-common-pair}
The pair $P_N(e)$ is controlled on every
$R\in\mathcal T(e)$ with $N\ge N_{\min}(R,e)$.
Up to its stop, its mass and vacuum sequence agree
with those on the top torus; its activity and
external scalar families agree there on flat polymers.
For every such level,
\begin{equation}\label{eq:rgfl-common-bounds}
 |m_j|\le\tfrac18e_j^{3/4-8\varepsilon},\qquad
 \|E_j\|_j\le\|\mathscr E_j\|_j^x
                             \le4e_j^{1/4-6\varepsilon},\qquad
 |\varepsilon_j^{\mathrm v}|
                         \le5c_{\mathrm{fl}}e_j^{1/4-6\varepsilon}.
\end{equation}
The real mass sequence is independent of the side
datum and of the twist. At all $N\ge0$,
\begin{equation}\label{eq:rgfl-bare-mass-bound}
 |m_{0,N}(e)|\le\tfrac18q^{3/4-8\varepsilon}<\tfrac18,
                       \qquad q=eL^{-N/2}.
\end{equation}
The total lattice mass is
$\boldsymbol\mu_N=\bar m\eta+m_{0,N}(e)$.
\end{theorem}

\begin{proof}
The pair is defined because
$N_{\min}(R_{\mathrm{top}},e)\le N_{\min}(R,e)$.
The two trajectories have the same initial
couplings and zero initial activities.
Suppose they agree through $j<k_*(R)$.
The top mass is bounded by $w_j^{\mathrm m}/8$;
the same bound on the smaller side permits continuation
to $j+1$ by \cref{prop:rgfl-continuation}.
Both inputs are now in the unit balls, are real,
and have locally agreeing activity and scalar families.
The scale proposition verifies their primitive,
domain and next-side comparison hypotheses.
\cref{thm:rgcov-step-comparison} gives the
same extracted mass and vacuum, including the same
$\delta\varepsilon_j^0$, and local agreement of the
outgoing activity and scalar family.
The scalar recursions therefore give the same
$m_{j+1},\varepsilon_{j+1}^{\mathrm v}$.
Restriction to zero external slot preserves agreement
of $E_{j+1}$.
This proves the induction up to $k_*(R)$.

The common mass bound follows from the top trajectory.
The local comparison alone is not an estimate of
the entire activity on the smaller torus; that
bound follows from its own continuation estimate.
The vacuum is a common scalar and satisfies the
top bound at this same level. Thus
\eqref{eq:rgfl-common-bounds} follows.
At a smaller-volume stop the couplings need not
vanish; the zeros defining $P_N(e)$ are imposed
at $k_{\mathrm{top}}$.

The fixed-point map on the top torus is determined
by $(N,e,\bar m)$, the fixed construction parameters
and photon coordinate basis, and the top primitive
families constructed from them.
The selected top side is itself a function of
$e,\bar m$ and these parameters.
The smaller side $R$ enters no part of this definition.
Nor does a twist: the twist is a harmonic part
of a connection inserted in the density identity,
whereas the coupling maps are maps of entire
local activity families and contain no such insertion.
The full forward activity recursion retains the external slot
as a variable, whereas the fixed-point map uses its restriction
to zero. No external connection or probe is inserted when
selecting the pair.
This proves the asserted independence within
the fixed construction.

For a cutoff on which the pair is defined,
\eqref{eq:rgfl-bare-mass-bound} is its initial mass
bound on the top torus. At the other finitely many
cutoffs it holds because the assigned mass is zero.
Finally $h\ge1$ implies $e\le\bar e_1\le1/3$,
so $q\le1/3$ and its positive power is below $1$.
The sequence is real in both cases.
The sign of $\boldsymbol\mu_N$ is not fixed by this
estimate; the lattice construction admits a real mass.
\end{proof}

On the reference torus, for $N\ge N_{\mathrm{ref}}(e)$,
the outputs needed at the terminal scale are
\begin{equation}\label{eq:rgfl-reference-terminal-input}
 a_*\in\{h/L,h/L^2\},\qquad
 |m_*|\le\tfrac18e_*^{3/4-8\varepsilon},\qquad
 \|E_*\|_*\le\|\mathscr E_*\|_*^x
                                 \le4e_*^{1/4-6\varepsilon}.
\end{equation}
The primitive estimates along this flow have the
same fixed constants on every torus and at every
cutoff. The remaining terminal density and partition
estimates concern the integration of these fields.



\begin{proof}[Proof of \Cref{thm:main}]
\Cref{prop:rgfl-scales} gives the stopping levels and thresholds;
\eqref{eq:rgfl-reference-stop} gives the reference terminal scales.
On the top torus, \cref{thm:rgfl-terminal-zero-pair} selects the unique
controlled pair with both terminal couplings zero.
\Cref{thm:rgfl-common-pair} transfers this pair and its scalar bounds
to each eligible smaller torus. Its proof gives the stated local
agreement of both activity families, while
\cref{prop:rgfl-continuation} supplies their full norm bounds.
The initial mass estimate and its zero extension are
\eqref{eq:rgfl-bare-mass-bound}. The same common-pair proof establishes
independence from the smaller side, twist and external insertion.
\end{proof}
\appendix
\section{The local formulas}\label{app:formulas}
This appendix recalls the finite expressions used in the comparison.
Their construction and estimates are the input results of
\cref{sec:inputs}. In the formulas below, $k$ is fixed and
$B=A+\mathscr C$.

\subsection{Extraction and reconstruction}
The extraction maps of \cite[\S3.1]{NormalizedContraction} use the
ordinary quadratic part
\[
 \mathscr E^{[2,0]}(X;A,\mathscr C)=\int\dd x\,\dd x'\,
 \bar\psi(x)^T K_X(A,\mathscr C;x,x')\psi(x').
\]
For a small $X$, choose one lift $\widetilde X$ and define
\begin{equation}\label{eq:rg-local-extractions}
 a_0(\mathscr E,X)
       =\operatorname{Vol}(X)^{-1}\mathscr E_{00}(X;0,0),\qquad
 a_2(\mathscr E,X)
       =\operatorname{Vol}(X)^{-1}
             \int_{\widetilde X}\dd x\int_{\widetilde X}\dd x'\,
                                                    K_X(0,0;x,x').
\end{equation}
Both cluster labels are averaged with measure $\eta_k^3/48$.
The second expression is independent of the lift: simultaneous
translation by a period multiplies its two arguments by the
same antiperiodic sign, whose product is one.

For an $M$-cube $\square$ put
\begin{equation}\label{eq:rg-counterterm-maps}
 \varepsilon_{\mathrm v}(\mathscr E,\square)
          =-\sum_{\substack{X\supset\square\\X\text{ small}}}
                                                     a_0(\mathscr E,X),
 \qquad
 M(\mathscr E,\square)
          =-\sum_{\substack{X\supset\square\\X\text{ small}}}
                                                     a_2(\mathscr E,X).
\end{equation}
Translations by $M\Z^3$ identify the small-polymer sets over
different cubes. They leave the two extracted quantities
unchanged. Thus the two sums do not depend on $\square$.
Write their values as $\varepsilon_{\mathrm v}(\mathscr E)$ and
$M(\mathscr E)$.

The imported scalarity assertion is $M(\mathscr E)=m(\mathscr E)1$.
On a small polymer the normalized remainder is
\cite[(17)]{NormalizedContraction}
\begin{equation}\label{eq:rg-normalized-remainder}
 (R\mathscr E)(X;A,\mathscr C)=\mathscr E(X;A,\mathscr C)
 -a_0(\mathscr E,X)\operatorname{Vol}(X)
 -\int_X\dd x\,\bar\psi(x)^Ta_2(\mathscr E,X)\psi(x).
\end{equation}
On other polymers it leaves the activity unchanged.
The operation $\Theta_{X,B}$ is the
shortest-chain quadratic adjustment of \cite[\S3.2]{NormalizedContraction}.
It averages over the shortest chains of $M$-cubes in $X$, using the
same six-path transports and pair coefficients $d_X^\alpha$.
We use the resulting $(R\mathscr E)^{\mathrm{nat}}$ from that definition.

Let $\Delta_y$ be the unit block centered at $y$.
For $c\in\Delta_y$, let $c_j(c)$ be the center of its block of side
$L^j\eta_k$. The averaging phase and its normalized average are
\[
 \varphi_A(y,c)=\exp\left(ie_k\eta_k\sum_{j=0}^{k-1}
          (\tau A)(c_{j+1}(c),c_j(c))\right),\qquad
 (\mathcal Q_k(A)f)(y)=\frac{\eta_k^3}{48}
       \sum_{c\in\Delta_y,g}\varphi_A(y,c)f(c,g).
\]
The reverse is the weighted transpose $\mathcal Q_k(-A)^T$.
For $y\in Y$ the two reconstruction rows of
\cite[Theorem~B]{NormalizedContraction} are as follows, where
$c_3=\pi x^{(3)}$.
\begin{equation}\label{eq:rg-reconstruction-two-pieces}
 \begin{split}
 (T_Y^{(1)}(A)\psi)(y)
 &=\int_{\Delta_y}\dd x\,\varphi_A(y,c)
    \Bigl\{(1+\mu/b_k)\psi(x)\\
 &\qquad-(b_k\eta_k)^{-1}\!\!\sum_{i:\,x^{(i)}\text{ over }Y}
           M_i(x)[U_A(x,x^{(i)})\psi(x^{(i)})-\psi(x)]\Bigr\},\\
 (T_Y^{(2)}(A)\chi)(y)
 &=-b_k^{-1}\eta_k^{\alpha-1}\int_{\Delta_y}\dd x\,
                1[c_3\notin Y]\varphi_A(y,c)M_3(x)\chi(x,x^{(3)}).
 \end{split}
\end{equation}
Rows outside $Y$ are zero. The first row contains ordinary
legs only over $Y$. The second row contains precisely the
type-$3$ pairs leaving $Y$, which are admissible boundary
pairs of length $\eta_k$. Their coefficients read $A$ only
on bonds with an endpoint in $\Delta_y$.
The barred rows are $\bar T_Y^{(i)}(B)=\mathsf C^{-1}T_Y^{(i)}(-B)\mathsf C$.
The substitution $\operatorname{sub}_{Y,B}$ replaces independent
unit generators $\Psi,\bar\Psi$ by these ordinary and pair rows and
fixes the old generators. This specifies the homomorphism used below.

\subsection{The determinant primitive and the first vertex}
Use the unprimed polymer resolvents $S_{k,y}(X,B)$ of
\cite[\S4.3]{FermionBlockManuscript}. Write
$\mathcal B_{k,y}=B_{k,y}$ for the finite block resolvent compared in
\cref{subsec:rgcov-finite-rows}, and $b_L=L^{-1}$.
At zero $F_0=D_B+\mu+b_LJP_LJ^*$ and $\Gamma_0=F_0^{-1}$.
The determinant primitive is \cite[(76)]{FermionBlockManuscript}:
\begin{align}
 \widetilde E_k(X,A)
    &=-ib_k^2\int_0^\infty\dd y\,
       \operatorname{Tr}\bigl[\gamma_5\mathcal B_{k,y}\mathcal Q_k
           S_{k,y}(X,A)\mathcal Q_k^\sharp\mathcal B_{k,y}\bigr]
                   \quad(k\ge1),\label{eq:local-determinant-definition}\\
 \widetilde E_0(X,A)
    &=\int_1^\infty\frac{\dd y}{y}
                     \operatorname{Tr}[F_0S_{0,y}(X,A)]
       -i\int_0^1\dd y\,\operatorname{Tr}[\gamma_5S_{0,y}(X,A)],\notag\\
 E_k^d(X,A)&=\tfrac14\bigl(\widetilde E_k(X,A)
       +\widetilde E_k(X,-A)
       +\overline{\widetilde E_k(X,\overline A)}
       +\overline{\widetilde E_k(X,-\overline A)}\bigr),\notag\\
 E_k^{\mathrm{det}}(X,A)&=E_k^d(X,A)-E_k^d(X,0).\notag
\end{align}
The traces use unit counting measure for coarse spinors and the
chamber measure $1/48$ at level zero. Absolute convergence and the
reblocked gain are part of the determinant input.

At zero write $q=e_0$ for the running charge. The local vertex of
\cite[(47)]{NonlinearStep} is
\begin{equation}\label{eq:rg-step-first-vertex}
 \begin{split}
 \mathcal V_0(Z,\square)=\frac12\int_\square\dd x\,
   \Big\{&
    \bar\psi(x)^TM_3(x)(e^{iqZ(c,c_3)}-1)
                                  [\psi(x)+\chi(x,x^{(3)})]\\
    &+[\bar\psi(x)+\bar\chi(x,x^{(3)})]^T
               M_3(x^{(3)})(e^{-iqZ(c,c_3)}-1)\psi(x)\Big\},
 \end{split}
\end{equation}
where $x=(c,g)$, $c_3=c+ge_1$ and
$\int_\square\dd x=48^{-1}\sum_{c\in\square}\sum_g$.
The pair $(x,x^{(3)})$ has $d_X=1$ and is admissible
for $\square$ of the second kind. Thus the vertex
has ordinary legs over the cube and precisely the
boundary pair legs allowed in the activity algebra.
Its coefficients are entire functions of $Z$ on
the bonds meeting the cube.


\subsection{The two finite localizations}
The definitions in this subsection are those of
\cite[\S5]{NonlinearStep}. For $\zeta\in\mathrm{Box}_k$ put
\begin{equation}\label{eq:rgnl-weakened-photon}
 Z(s')=H_k^{\mathrm{ph}}(s')C C_k^{1/2}(s')\zeta.
\end{equation}
The radii, with $e=e_k$, are
\begin{equation}\label{eq:rgnl-contour-radii}
 r_1=\tfrac12e^{-3/4+9\varepsilon/2},\qquad
 r_2=e^{-3/4+7\varepsilon},\qquad r_4=e^{-1/4+3\varepsilon}.
\end{equation}
The imported contour bounds give $2\le r_2,r_4\le r_1$ and
$A+tZ\in\widetilde R_k$, $B+tZ\in\widetilde G_k$ for $|t|\le r_1$.
The arguments of $H_k^\#$ are independent pairs
$(\Phi,\bar\Phi)$; in $H_k^\#\Phi$ the single symbol denotes
this pair. Write
$\psi^\#=(\psi,\chi;\bar\psi,\bar\chi)$ for the ordinary,
pair, barred ordinary and barred pair generators, respectively.
Then
\[
 \begin{aligned}
 H_k^\#(s'',B)\Phi
   ={}&(H_k(s'',B)\Phi,\delta_{\alpha,B}H_k(s'',B)\Phi;\\
       &
       \bar H_k(s'',B)\bar\Phi,
                       \delta_{\alpha,-B}\bar H_k(s'',B)\bar\Phi),\\
 \bar H_k(s'',B)&=\mathsf C^{-1}H_k(s'',-B)\mathsf C.
 \end{aligned}
\]
Define
$J_k^\#(s'',B,Z)=H_k^\#(s'',B+Z)-H_k^\#(s'',B)$ and
\[
 \widetilde\Gamma_k(s'',B)(W,\bar W)
                         =(\Gamma_k(s'',B)W,\bar W),\qquad
 W_k^\#(s'',B)=H_k^\#(s'',B)
                         \big(\Gamma_k(s'',B)W,\bar W\big).
\]
The notation $J_k,\bar J_k$ denotes the components without
H\"older quotient; $\mathrm{Wc},\overline{\mathrm{Wc}}$ denotes
the corresponding components of $W_k^\#$.
The auxiliary generators $W$ have weight
$\mathsf w_W=1$ for $k\ge1$ and $\sqrt{48}$ for $k=0$.

In the following localization, cubes and polymers are sets
of photon sites. Containment refers to all their lattice sites,
and two cubes meet when they share a lattice site.
Let $P_0(X)$ be the parameter cubes contained in $X$.
For a set $P$ of parameter cubes disjoint from $P_0(X)$, let $Y'(X,P)$ be the union of
the $M$-cubes meeting $X$ or a cube of $P_0(X)\cup P$.
For a continuous analytic family $G(s,X)$ define
\begin{align}
 \mathcal T_PG(X)
  &=\int_{[0,1]^P}\dd s_P\,
       \partial_PG(s,X)
       \big|_{s=1\text{ on }P_0(X),\,
                    s=0\text{ outside }P_0(X)\cup P},
                          \label{eq:rgnl-interpolation}\\
 \operatorname{Loc}G(Y)
  &=\sum_{(X,P):Y'(X,P)=Y}\mathcal T_PG(X).
                          \label{eq:rgnl-localization}
\end{align}
For a function analytic near the closed disk $|t|\le r$, $r>1$,
write
\[
 \mathcal C_rf=\frac1{2\pi i}\int_{|t|=r}
                                      \frac{f(t)}{t(t-1)}\,\dd t .
\]
This is $f(1)-f(0)$ and satisfies
$\|\mathcal C_rf\|\le(r-1)^{-1}\sup_{|t|=r}\|f(t)\|$.
Let $(m,\mathscr E)\in\operatorname{Ball}_k^x(4)$ and set
$m'=m+m(\mathscr E)$, $\mathscr E^r=R\mathscr E$.
For $k\ge1$ define
\begin{align}
 G_1(s,X)&=\mathcal C_{r_1}
                \langle\bar\Psi,M_k(X,B+tZ(s'))\Psi\rangle_1,
                                      \label{eq:rgnl-G1}\\
 G_2(s,X)&=\mathcal C_{r_2}\,
       \mathscr E^r(X;A+tZ(s'),\mathscr C,
                    \psi^\#+tJ_k^\#(s'',B,Z(s'))\Psi),
                                      \label{eq:rgnl-G2}\\
 G_3(s,\square)&=\mathcal C_{r_2}(-m')
       \langle\bar\psi+t\bar J_k(s'',B,Z(s'))\bar\Psi,
                         \psi+tJ_k(s'',B,Z(s'))\Psi\rangle_\square .
                                      \label{eq:rgnl-G3}
\end{align}
The symbols $M_k(X,\cdot)$ denote the local Dirac primitives.
A function defined for one cube is zero on the other polymers.
Put
\[
 E_i^*(Y)=\operatorname{Loc}G_i(Y),\qquad
 E_i(Y)=\operatorname{sub}_{Y,B}E_i^*(Y),\qquad
 E_{\le3}=\sum_{i=1}^3E_i.
\]
The substitution was defined above.
At $k=0$ there is no $\Psi$ and no substitution:
$G_1(s,\square)=\mathcal V_0(Z(s'),\square)$,
$G_2=G_3=0$, and $E_{\le3}=\operatorname{Loc}G_1$.
All parameter kinds occur in this first localization.

For the second localization, which uses only $\pi_k$, let
$v=W_k^\#(s'',B)$ and define
\begin{align}
 G_{\le3}^{\,W}(s'',X)&=E_{\le3}(X;A,\mathscr C,\zeta,
                                                      \psi^\#+v),
                                      \label{eq:rgnl-GW}\\
 G_4(s'',X)&=\mathcal C_{r_4}
                  \mathscr E^r(X;A,\mathscr C,\psi^\#+tv),\\
 G_5(s'',\square)&=\mathcal C_{r_4}(-m')
                 \langle\bar\psi+t\overline{\mathrm{Wc}},
                                   \psi+t\mathrm{Wc}\rangle_\square,\\
 F(Y)&=\operatorname{Loc}(G_{\le3}^{\,W}+G_4+G_5)(Y).
                                      \label{eq:rgnl-finite-activity}
\end{align}
At $k=0$ the term $G_4$ is zero and $m'=m$.


\subsection{Integration, logarithm and reblocking}
For the even activity $F$ just defined, set
\begin{align*}
 K(Y)&=\sum_{\substack{\{X_i\}\text{ distinct}\\\bigcup_iX_i=Y}}
             \prod_i(\exp F(X_i)-1),\qquad K(\varnothing)=1,\\
 H(Y)&=\int\dd\mu_I(W)\,\dd\widehat\mu_k(\zeta)\,K(Y),\\
 E_k^\#(Y)&=\sum_{n\ge1}\frac1{n!}
       \sum_{\bigcup_iY_i=Y}u^T(Y_1,\ldots,Y_n)\prod_iH(Y_i).
\end{align*}
Here $\widehat\mu_k$ is the product of the normal laws restricted
and normalized on $[-p_{0,k},p_{0,k}]$. The measure $\mu_I(W)$ is
the normalized identity-covariance Grassmann measure: the covariance
kernel is the inverse mode weight. The coefficient $u^T$ is the
sum of $(-1)^{|E(G)|}$ over connected spanning subgraphs of the
incompatibility graph of the tuple. Two polymers are incompatible
when their union is connected; in particular repetitions are incompatible.
Absolute convergence in the weighted coefficient norm is supplied by
\cite[Proposition~6.1]{NonlinearStep}.

For an outgoing polymer $Y$, put
\[
 E_k^*(Y;A',\mathscr C')=\sum_{\bar X^L=LY}
              S_LE_k^\#(X;A'_L,\mathscr C'_L),\qquad
 A'_L(b)=L^{-1/2}A'(b/L).
\]
The bar denotes the union of $LM$-cubes containing the $M$-cubes of $X$.
The map $S_L$ multiplies ordinary and pair legs by $L^{-1}$ and
$L^{-1-\alpha}$, respectively. The same sum defines $\mathcal L$
on any input activity. The scaled-domain inclusion used in comparison is
\begin{equation}\label{eq:rg-scaled-domain-inclusion}
 (A',\mathscr C')\in\widetilde R_{k+1}\times\mathfrak X_{k+1}
 \quad\Longrightarrow\quad
 (A'_L,\mathscr C'_L)\in\widetilde R_k^-\times\mathfrak X_k;
\end{equation}
it is \cite[(29)]{NormalizedContraction} under the assumed numerical conditions.
\printbibliography
\end{document}
